\documentclass[11pt]{amsart}
%\usepackage[margin=1in]{geometry}  % set the margins to 1in on all sides
\usepackage{graphicx}              % to include figures
\usepackage{amsmath}              % great math stuff
\usepackage{amsfonts}              % for blackboard bold, etc
\usepackage{amsthm}                % better theorem environments
%\usepackage{showkeys}
\usepackage{color}
\usepackage{tabulary}
%\spackage{slashbox}
\usepackage{caption}
\usepackage{subcaption}
\usepackage{amssymb}
\usepackage{todonotes}
\theoremstyle{plain}
\usepackage{mathtools}
%\usepackage{showkeys}
%\usepackage{showlabels}
% various theorems, numbered by section

\theoremstyle{plain}
\newtheorem{thm}{Theorem}
\newtheorem*{cor}{Corollary}
\newtheorem{theorem}{Theorem}[section]
\newtheorem{question}[theorem]{Question}
\newtheorem{proposition}[theorem]{Proposition}
\newtheorem{lemma}[theorem]{Lemma}
\newtheorem{claim}[theorem]{Claim}
\newtheorem{fact}[theorem]{Fact}
\newtheorem{observation}[theorem]{Observation}
\newtheorem{conjecture}[theorem]{Conjecture}
\newtheorem{corollary}[theorem]{Corollary}
\theoremstyle{definition}
\newtheorem{definition}[theorem]{Definition}
\newtheorem{remark}[theorem]{Remark}
\newtheorem{example}[theorem]{Example}

\newcommand{\nc}{\newcommand}
\nc{\dmo}{\DeclareMathOperator}

\nc{\Q}{\mathbb{Q}}
\nc{\R}{\mathbb{R}}
\nc{\Z}{\mathbb{Z}}
\nc{\N}{\mathbb{N}}
\nc{\C}{\mathbb{C}}
\nc{\cS}{\mathcal{S}}
\nc{\iso}{\cong}
\dmo{\Mod}{Mod}
\dmo{\Diff}{Diff}
\dmo{\Homeo}{Homeo}
\dmo{\dist}{dist}
\dmo\BDiff{BDiff}
\dmo\SO{SO}
\dmo\slide{sl}
\dmo\im{im}
\dmo\id{id}
\dmo\Fix{Fix}
\dmo\Out{Out}
\dmo{\T}{\mathcal{T}}
\dmo{\Te}{\mathcal{T}^{\epsilon}}
\dmo{\M}{\mathcal{M}}
\dmo{\Me}{\mathcal{M}^{\epsilon}}

\renewcommand{\epsilon}{\varepsilon}
\nc{\coloneq}{\mathrel{\mathop:}\mkern-1.2mu=}
\nc{\margin}[1]{\marginpar{\scriptsize #1}}
\nc{\para}[1]{\bigskip\noindent\textbf{#1}}

\DeclarePairedDelimiter\ceil{\lceil}{\rceil}
\DeclarePairedDelimiter\floor{\lfloor}{\rfloor}




\newcommand{\bd}[1]{\mathbf{#1}}  % for bolding symbols

      
\newcommand{\col}[1]{\left[\begin{matrix} #1 \end{matrix} \right]}
\newcommand{\comb}[2]{\binom{#1^2 + #2^2}{#1+#2}}
\newcommand{\into}{\hookrightarrow}
\newcommand{\onto}{\twoheadrightarrow}
\newcommand{\ie}{i.e.\ }
\newcommand{\mathcalo}{\mathcal{O}}
\newcommand{\inj}{inj }
\newcommand{\vol}{vol }
\newcommand{\cal}{\mathcal}

\begin{document}


\title[Small eigenvalues and thick-thin decomposition]
{Small eigenvalues and thick-thin decomposition in negative curvature}
\author{Ursula Hamenst\"{a}dt}

\thanks
{AMS subject classification: 58J50,53C20\\
Keywords: Spectre du Laplacien, conditions de bord Neumann,
vari\'et\'es de courbure n\'egative pinc\'ee, d\'ecomposition
en parties \'epaisse et fine. \\
Spectrum of the Laplacian, Neumann boundary conditions, 
manifolds of pinched negative curvature, 
thick-thin decomposition\\
Research
supported by ERC grant ``Moduli''}
%\\
%e-mail: ursula@math.uni-bonn.de}
\date{November 14, 2019}


\begin{abstract}
Soit $M$ une vari\'et\'e Riemannienne compl\`{e}te orient\'ee,
de dimension $n\geq 3$ et de volume finie. Supposons que la
courbure de $M$ soit contenue dans $[-b^2,-1]$, et soit
$M=M_{\rm thick}\cup M_{\rm thin}$ la d\'ecomposition 
en sa partie \'epaisse et sa partie fine. Soit $\lambda_k(M)$ la 
$k$-ti\`{e}me valeur
propre de l'op\'erateur Laplacien, avec conditions de bord de Neumann. 
Nous d\'emontrons que 
$\lambda_k(M_{\rm thick})/3\leq \lambda_k(M)$ pour tout $k$ tel que
$\lambda_k(M)<(n-2)^2/12$. 
Si $M$ est hyperbolique et de dimension 3, alors $\lambda_k(M)\leq 
C\log({\rm vol}(M_{\rm thin})+2)\lambda_k(M_{\rm thick})$ pour un nombre
$C>0$ fix\'e pourvu que
$\lambda_k(M_{\rm thick})<1/96$.

\bigskip

Let $M$ be a finite volume oriented complete Riemannian
manifold of dimension $n\geq 3$ and curvature in $[-b^2,-1]$, 
with thick-thin decomposition $M=M_{\rm thick}\cup M_{\rm thin}$.
Denote by $\lambda_k(M_{\rm thick})$ 
the $k$-th eigenvalue for the Laplacian on $M_{\rm thick}$, with 
Neumann boundary conditions. We show  
that $\lambda_k(M_{\rm thick})/3\leq \lambda_k(M)$ for all $k$ for which 
$\lambda_k(M)<(n-2)^2/12$. 
If $M$ is hyperbolic and of dimension 3 then $\lambda_k(M)\leq 
C\log({\rm vol}(M_{\rm thin})+2)\lambda_k(M_{\rm thick})$ for a fixed
number $C>0$ 
provided that $\lambda_k(M_{\rm thick})<1/96$.
%The first inequality is true in any dimension.
%and that this estimate is roughly sharp.
\end{abstract}

 

\maketitle

%\tableofcontents

\section{Introduction}


The small part of the spectrum of the Laplace operator $\Delta$ acting on 
functions on a closed oriented 
hyperbolic surface $S$ is quite well understood. Namely,
if $g$ denotes the genus of $S$ then the $2g-3$-th eigenvalue 
$\lambda_{2g-3}(S)$ can be arbitrarily
small \cite{Bu}, while $\lambda_{2g-2}(S)>\frac{1}{4}$ \cite{OR09}.

By the Gauss-Bonnet theorem, the volume of $S$ equals $2\pi(2g-2)$, so
these results relate the small part of the spectrum of $S$ to its volume.

For $n\geq 3$, the small part of the spectrum of an oriented  
finite volume Riemannian 
manifold $M$ of dimension $n$ and 
sectional curvature $\kappa\in [-b^2,-1|$ for some $b\geq 1$ 
is less well understood. 
The manifold $M$ admits a \emph{thick-thin decomposition}
$M=M_{\rm thick}\cup M_{\rm thin}$ 
which is determined as follows. There exists a constant
$c=c(n)>0$ such that $\epsilon=b^{-1}c(n)$ is a \emph{Margulis constant} 
for $M$, and 
$M_{\rm thick}$ is the set of all points in $M$ of injectivity radius 
at least $\epsilon$. 
Its complement $M_{\rm thin}$ is a disjoint union of
so-called \emph{Margulis tubes} and \emph{cusps}. 
Here a Margulis tube is a tubular neighborhood of a closed geodesic 
of length at most $2\epsilon$, and a cusp is homeomorphic to the quotient of 
a horoball in the universal covering $\tilde M$ of $M$ 
by a rank $n-1$ parabolic subgroup of the isometry group of $\tilde M$. 
After a small modification, 
we may assume that
$M_{\rm thick}=M-M_{\rm thin}$ is a submanifold of $M$ with smooth boundary
\cite{BCD93}.

Unlike in the case of surfaces, 
the submanifold $M_{\rm thick}$ is always connected, and
this prevents the occurence of very small eigenvalues. 
More precisely, for every $n\geq 3$, 
Schoen \cite{S82} established the existence of a universal and explicit 
constant $\theta=\theta(n,b)>0$ such that
\begin{equation}\label{lower}\lambda_1(M)\geq 
\frac{\theta}{{\rm vol}(M)^2}\end{equation}
for \emph{every} closed $n$-manifold $M$ with curvature in 
$[-b^2,-1]$ (here in contrast to the work of Schoen, we 
normalize the metric on $M$ so that the upper bound 
of the curvature is fixed). That 
this estimate extends without
change to non-compact finite volume manifolds
was established in \cite{DR86,D87}.

For hyperbolic manifolds of dimension $n=3$, this bound is roughly sharp:  
White \cite{White} 
proved that for fixed $r \geq 2, \delta > 0$ 
there exists a constant 
 $a=a(r,\delta)>0$ such that 
 \[ \lambda_1(M)\in [1/a {\rm vol}(M)^2, a/{\rm vol}(M)^2]\]
 for any closed 
hyperbolic 3-manifold whose injectivity radius is bounded from 
below by $\delta$ and such 
that the rank of the fundamental group $\pi_1(M)$ of $M$ is at most $r$.
This rank is defined to be the minimal number of generators of 
$\pi_1(M)$, and it 
is bounded from above by a fixed multiple of the
volume (Theorem 1.10 in \cite{G04}). A similar statement also holds
true for random 3-dimensional hyperbolic mapping tori \cite{BGH}
and for random hyperbolic 3-manifolds of fixed Heegaard genus
\cite{HV19}. 



Since the injectivity radius at points in $M_{\rm thick}$ is 
at least $\epsilon$, the submanifold $M_{\rm thick}$ of $M$ is 
uniformly
quasi-isometric to a finite graph $G$ of uniformly 
bounded valence, with constant
only depending on the dimension and the curvature bounds 
(see \cite{BGLM02} for a detailed discussion of this fact 
in the case of hyperbolic manifolds). 
If $\vert G\vert $ denotes 
the number of vertices of $G$, 
then for $k\leq \vert G\vert $ the $k$-th eigenvalue
$\lambda_k(M_{\rm thick})$ of $M_{\rm thick}$ with Neumann boundary
conditions is uniformly comparable to the $k$-th eigenvalue 
$\lambda_k(G)$ 
of the graph 
Laplacian of $G$ \cite{Man05}. 
Here and in the sequel, we list eigenvalues as 
$0=\lambda_0<\lambda_1\leq \lambda_2\leq \dots$
with each eigenvalue repeated according to its multiplicity. 
Note that $\vert G\vert $ is proportional to the 
volume ${\rm vol}(M_{\rm thick})$ of $M_{\rm thick}$, 
with multiplicative constants only depending on $n$ and $b$.

Now for any graph $G$, there are precisely $\vert G\vert$ eigenvalues
$0=\lambda_0(G)<\lambda_1(G)\leq\dots\leq \lambda_{\vert G\vert -1}(G)$, and
(Lemma 1 of \cite{Chu96})  
\[\sum_i\lambda_i(G)=\vert G\vert.\]
In particular, we have $\lambda_{\vert G\vert-1}(G)\geq 1$.
As the volume of a finite volume oriented manifold $M$ of curvature in 
$[-b^2,-1]$ is uniformly comparable to the volume of $M_{\rm thick}$, this 
together with the result in \cite{Man05} implies 
that there exists a constant
$q=q(n,b)>1$ so that $\lambda_{q{\rm vol}(M)}(M_{\rm thick})\geq 1/q$. 

To recover a relation between the spectrum of $M$ and the
volume of $M$ which resembles the result known for hyperbolic surfaces,
it is therfore  
desirable
to relate the small part of the spectrum of $M$ to the small part
of the spectrum of $M_{\rm thick}$, taken with Neumann boundary condition.
The first purpose of this
article is to establish such a relation. We show




\begin{thm} \label{thm1}
For a suitable choice of a Margulis constant, we have 
\[\lambda_k(M)\geq \min\{\frac{1}{3}\lambda_k(M_{\rm thick}),\frac{(n-2)^2}{12}\}\]
for every finite volume oriented Riemannian manifold $M$ of dimension
$n\geq 3$ and curvature 
$\kappa\in [-b^2,-1]$, and all $k\geq 0$. 
\end{thm}


The dependence of this estimate 
on the lower curvature bound $-b^2$
enters this result via the Margulis constant which depends on $b$.
Note also that the constant $(n-2)^2/12$ is 
uniformly comparable to
the lower bound $(n-1)^2/4$ for  
the bottom of the essential spectrum of a non-compact
finite volume manifold of curvature $\kappa\in [-b^2,-1]$
(Corollary 3.2 of \cite{H04}). 
We expect that our methods can be used to extend
Theorem \ref{thm1} 
to geometrically finite manifolds of infinite volume and
curvature in $[-b^2,-1]$. This could then be used to establish
an improvement of Corollary 3.3 of \cite{H04} which contains
the following statement as a special case: 
The number of eigenvalues
contained in $(0,(n-2)^2/12)$ is at most $d^{{\rm vol}(M)}$ for a fixed 
constant $d>0$. However, we do not attempt to carry 
out such a generalization in this article. 

As an application of Theorem \ref{thm1}, we recover the results of Schoen, 
of Dodziuk and of Randol and Dodziuk, and we relate 
the number of small 
eigenvalues to the volume 
as promised. 

\begin{cor}\label{application}  
For all $n\geq 3,b\geq 1$ there exists a constant 
$\chi=\chi(n,b)>0$ with the
following property. 
Let $M$ be a finite volume oriented Riemannian 
manifold of dimension $n$ and curvature $\kappa\in [-b^2,-1]$; then 
\begin{enumerate}
\item
$\lambda_1(M)\geq \frac{\chi}{{\rm vol}(M)^2}.$
\item $\lambda_{ {\rm vol}(M)/\chi}(M)\geq \chi$.
\end{enumerate}
\end{cor}




Unlike in the work of Schoen, the constant $\chi(n,b)$ 
in the above Corollary 
is not explicit as it 
depends on a Margulis constant for Riemannian 
manifolds with curvature in $[-b^2,-1]$.
However, for hyperbolic manifolds
it can explicitly be estimated. 

For hyperbolic 3-manifolds $M$ we also obtain upper bounds  
for the small eigenvalues of $M$. We show

\begin{thm}\label{thm2}
There exists a number $c>0$ such that for every finite volume 
oriented hyperbolic 3-manifold $M$, we have  
\[\lambda_k(M)\leq c\log ({\rm vol}(M_{\rm thin})+3)\lambda_k(M_{\rm thick})\]
for all $k\geq 1$ such  that $\lambda_k(M_{\rm thick})<1/96$.
\end{thm}

%Our second result shows that this estimate is sharp.
%If the injectivity radius of $M$ is bigger than $\epsilon$ then clearly
%$\lambda_1(M)=\lambda_1(M_{\rm thick})$, so the lower bound for
%$\lambda_1(M)$ is clearly sharp. We show that the upper bound is 
%also sharp.

%\begin{thm}\label{thm2}
%There exists a constant $c^\prime>0$ and a sequence $M_i$ of 
%closed hyperbolic 3-manifolds with 
%${\rm vol}(M_i)\to \infty$ and
%\[\lambda_1(M)\geq c^\prime \log{\rm vol}(M_{\rm thin})\lambda_1(M_{\rm thick}).\]
%\end{thm}

The proof of Theorem \ref{thm1} uses standard comparison results and a
simple decomposition principle, and it is 
carried out in Section 2. Theorem \ref{thm2}
is shown in Section \ref{above}
with an explicit construction which is only 
valid for hyperbolic 3-manifolds. 


\bigskip\noindent 
{\bf Acknowledgement:} I am grateful to Juan Souto for
inspiring discussions. Some versions of the results in this paper
were independently obtained by 
Anna Lenzhen and Juan Souto  \cite{LS16}.
%Given the discussions we had, we felt it would be most
%appropriate if I added Anna Lenzhen and Juan Souto 
%as authors to this paper. 
%I very much regret that they 
%declined my invitation.

 



%%%%%%%%%%%%%%%%%%%%%%%%%%%%%%%%%%%%%%%%%%%%%%%%%%%%%%%%%%%%%%%%%%%%%%%%%%%
%%%%%%%%%%%%%%%%%%%%%%%%%%%%%%%%%%%%%%%%%%%%%%%%%%%%%%%%%%%%%%%%%%%%%%%%%%%



\section{Bounding small eigenvalues from below}\label{below}

The goal of this section is to show Theorem \ref{thm1}.

Thus let $M$ be a finite volume oriented Riemannian 
manifold of dimension $n\geq 3$
and sectional curvature $\kappa\in [-b^2,-1]$ for some 
$b\geq 1$.  
Then $M$ admits a \emph{thick-thin decomposition}
\[M=M_{\rm thin}\cup M_{\rm thick}.\]
For a number $\epsilon=b^{-1}c(n)>0$, a so-called \emph{Margulis constant}, 
the thin part $M_{\rm thin}$ is the set of all 
points $x\in M$ with injectivity
radius ${\rm inj}(x)\leq  \epsilon$,
and $M_{\rm thick}=\{x\mid {\rm inj}(x)\geq \epsilon\}$.
The set 
$M_{\rm thick}$ is a non-empty 
compact connected manifold with (perhaps non-smooth) boundary, and  
$M_{\rm thin}$ is a union of (at most) finitely many 
\emph{Margulis tubes} and \emph{cusps}. 


A Margulis tube is a tubular neighborhood of a closed geodesic
$\gamma$ in $M$ of length smaller than $2\epsilon$, and it is 
homeomorphic to $B^{n-1}\times S^1$
where $B^{n-1}$ is the closed unit ball in $\mathbb{R}^{n-1}$.
The geodesic $\gamma$ 
is called the \emph{core geodesic} of the tube. 

Let $T$ be such a Margulis tube, with core geodesic $\gamma$
of length $\ell<2\epsilon$. 
We fix a parameterization of $\gamma$ by arc length on the interval
$[0,\ell)$. 
Let $\sigma$ be the standard angular coordinate on the fibers of the 
unit normal bundle $N(\gamma)$ 
of $\gamma$ in $M$ obtained by parallel transport of 
the fibre over $\gamma(0)$  
(this unit normal bundle is an $S^{n-2}$-bundle
over $\gamma$), let $s$ be the length parameter
of $\gamma$ and let $\rho\geq 0$ be the radial distance from $\gamma$.
Via the normal exponential map, these functions define 
"coordinates" (i.e. a parameterization) $(\sigma,s,\rho)$ for $T-\{\gamma\}$,
defined on an open subset of $N(\gamma)\times (0,\infty)$
which will be specified below. 
In these "coordinates",
the maps $\rho\to (\sigma, s, \rho)$ are unit speed geodesics 
with starting point on $\gamma$ and initial velocity
perpendicular to $\gamma^\prime$.

There exists a continuous function 
\[R:N(\gamma)\to (0,\infty), (\sigma,s)\to R(\sigma,s)\]
such that in these "coordinates", we have
$T=\{\rho\leq R\}$. The metric on $T-\{\gamma\}$ is 
of the form
$h(\rho) +d\rho^2$ where $h(\rho)$ 
is a family of smooth metrics on the hypersurfaces $\rho={\rm const}$.

Lemma 2.4 of \cite{BCD93} states that there exists a constant
$\nu(n,b)>0$ which can be computed as an explicit function 
of the constants $c(n), b,n$ such that
\begin{equation}\label{mincontrol}
\min R\geq -\log \ell -\nu(n,b).\end{equation}
In particular, up to slightly adjusting the thick-thin decomposition
and replacing $M_{\rm thick}$  
by its union with all Margulis tubes with core geodesics of length
$\ell$ so that $\log \ell\geq -3-\nu(n,b)$, 
we may assume that 
for every component $T$ of $M_{\rm thin}$, the distance between
the core geodesic and the boundary $\partial T$ is at least three. 

In general, the boundary of a Margulis tube need not be smooth.
However, 
Theorem 2.14 of \cite{BCD93} shows that it can be perturbed 
to be smooth and of controlled geometry. 
We record this result for completeness.

\begin{theorem}\label{thmbcd93}
Let $T\subset M$ be a Margulis tube, with core geodesic $\gamma$ and 
boundary $\partial T$.
Then there exists a smooth hypersurface $H\subset T-\gamma$ with
the following properties.
\begin{enumerate}
\item The angle $\theta$ between the tangent of the radial geodesic
and the exterior normal to $H$ is less than $\pi/2-\alpha$
for some $\alpha=\alpha(n,b)\in (0,\pi/2)$.
\item The sectional curvatures of $H$ with respect to the induced
metric are bounded in absolute value by a constant depending only
on $n$ and $b$.
\item $H$ is homeomorphic to $\partial T$ by pushing along radial
arcs. The distance between $x\in H$ and its image 
$\bar x\in \partial T$ satisfies $d(x,\bar x)\leq bc(n)/50$.
\end{enumerate} 
\end{theorem}

In the sequel we always assume that the boundary $\partial T$ of 
a Margulis tube has the properties stated in 
Theorem \ref{thmbcd93}.
Then the injectivity radius of $\partial T$ 
with respect to the induced metric is bounded from below by
a positive constant only depending on the curvature bound and
the dimension (Corollary 2.24 of \cite{BCD93}). 


Our first goal is to obtain a better understanding of the 
volume element of a Margulis tube. To this end 
we record a variant of Lemma 1 
of \cite{DR86}. 
We begin with a simple comparison lemma. 

Let $\tilde M$ be a simply connected complete Riemannian manifold 
of dimension $n$, with sectional curvature
$\kappa\leq -1$. Let $\eta:\mathbb{R}\to \tilde M$ be a geodesic parametrized by 
arc length and let $Y_1(t),\dots,Y_{n-1}(t)$ be orthonormal parallel
vector fields along $\eta$, orthogonal to $\eta^\prime$. Let moreover 
$J_1,\dots,J_{n-1}$ be Jacobi fields along $\eta$,
orthogonal to $\eta^\prime$, with the following initial conditions.

\begin{enumerate}
\item $J_1(0)=Y_1(0)$, and 
$\frac{\nabla}{dt}J_1(0)=0$.
\item For $i=2,\dots,n-1$ 
we have  
$J_i(0)=0$, and covariant derivatives 
$\frac{\nabla}{dt} J_i(t)\vert_{ t=0}=Y_i(0)$. 
\end{enumerate}

Then $J(t)=(J_1(t),\dots,J_{n-1}(t))$ can be viewed as an 
$(n-1,n-1)$-matrix with respect to the basis $Y_1(t),\dots,Y_{n-1}(t)$.  
Define the matrix $A(t)$ by 
\begin{equation}\label{riccati}
\frac{\nabla}{dt}J(t)=A(t)J(t).\end{equation}
The matrix valued map $t\to A(t)$ satisfies the \emph{Riccati equation}
\[A^\prime +A^2+R_{\eta^\prime}=0\]
where $R_{\eta^\prime}$ is the curvature tensor of $\tilde M$ 
evaluated on $\eta^\prime$. In particular,
$A$ is self-adjoint. 




Using the notations from \cite{E94}, 
for $t>0$ we denote by $j(t)$ the determinant of the matrix $J(t)$. 


\begin{lemma}\label{volumegrowth}
For all $R>0$ we have 
\[j(R)\geq (n-1)\int_0^Rj(t)dt.\]
\end{lemma}
\begin{proof}
The proof follows from standard comparison. 
We use the above notations. 

Let $\bar \eta:\mathbb{R}\to \mathbb{H}^n$ be a geodesic in the hyperbolic $n$-space $\mathbb{H}^n$.
Let $\bar Y_1(t),\dots,\bar Y_{n-1}(t)$ be parallel vector fields along $\bar\eta$
such that for each $t$, $\bar Y_1(t),\dots,\bar Y_{n-1}(t)$ is an orthonormal basis of 
$\bar \eta^\prime(t)^\perp$. 
Let $\bar J_i$ $(1\leq i\leq n-1)$ be the Jacobi fields along $\bar\eta$ defined by 
$\bar J_1(0)=\bar Y_1(0),\frac{\nabla}{dt}\bar J_1(0)=0$, 
and for $i\geq 2$ we require that
$\bar J_i(0)=0$ and $\frac{\nabla}{dt}\bar J_i(0)=\bar Y_i(0)$. Write 
$\bar J(t)=(\bar J_1(t),\dots,\bar J_{n-1}(t))$ and view this as a matrix with respect to the
basis $\bar Y_1(t),\dots,\bar Y_{n-1}(t)$ 
of $\bar\eta^\prime(t)^\perp$.


The Jacobi fields $\bar J_i$ can explicitly be computed as follows. We have
$\bar J_1(t)=\cosh(t)\bar Y_1(t)$, and 
$\bar J_i(t)=\sinh(t)\bar Y_i(t)$ for $i\geq 2$. 
In particular, if we denote by $\bar j(t)$ the determinant
of $\bar J(t)$ then 
\[\bar j(t)=\sinh^{n-2}(t)\cosh(t).\]
Thus $\bar j^\prime(t)=(n-2)\sinh^{n-3}(t)\cosh^2(t)+
\sinh^{n-1}(t)$ and hence
\[\bar j^\prime(t)=((n-2)\coth(t)+\tanh(t))\bar j(t).\] 


Using the notations from the lemma and the text preceding it, 
Theorem 3.2, Theorem 3.4 and Section 6.1 of \cite{E87} show that 
\begin{equation}\label{eschenburg}
j^\prime(t)/j(t)\geq \bar j^\prime(t)/\bar j(t)\end{equation}
for all $t>0$. 


Write $b(t)=\frac{1}{n-1}\sinh^{n-1}(t)$; then 
$b(t)=\int_0^t\bar j(s)ds$ and 
\begin{equation}\label{hyperbolic}
\frac{d}{dt}\log b(t)=\frac{b^\prime(t)}{b(t)}=(n-1)\coth(t)>n-1
\end{equation}
for all $t$. 

For $a(t)=\int_0^tj(s)ds$ we have $a^\prime(t)=j(t)$. 
By the estimate (\ref{hyperbolic}), it now suffices to show 
that 
\[\frac{d}{dt}\log a(t)=\frac{a^\prime(t)}{a(t)}\geq \frac{d}{dt}\log b(t)\]
for all $t>0$. 

Let $\delta>0$. It suffices to show that 
for all $t>0$ we have 
$\frac{d}{dt}\log a(t)\geq (1-\delta)\frac{d}{dt}\log b(t)$. 
To this end note that by comparison, the inequality holds  true for all small $t$ (see \cite{E87} for 
details). 
As this condition is closed, if the inequality does not hold for all $t$ then 
there is a number $T_0>0$ so that the 
estimate holds true for $t\leq T_0$, and it is violated for 
$T_0<t<T_0+\tau$ where $\tau >0$. 
Then we have 
$\frac{d}{dt}\log a(T_0)=(1-\delta)\frac{d}{dt}\log b(T_0)$, 


Now 
\[\frac{d^2}{dt^2}\log a=\frac{d}{dt}\frac{a^\prime}{a}=\frac{a^{\prime\prime}a-(a^\prime)^2}{a^2}
=\frac{a^{\prime\prime}}{a}-(\frac{a^\prime}{a})^2=\frac{a^\prime}{a}(\frac{a^{\prime\prime}}{a^\prime}-
\frac{a^\prime}{a}).\]
Inequality (\ref{eschenburg}) implies that 
$\frac{a^{\prime\prime}}{a^\prime}(T_0)\geq 
\frac{b^{\prime\prime}}{b^\prime}(T_0)$.
As $\frac{a^\prime}{a}(T_0)=(1-\delta)\frac{b^\prime}{b}(T_0)$,
we conclude that
\[\frac{d^2}{dt^2}\log a(T_0)> (1-\delta)
\frac{d^2}{dt^2}\log b(T_0).\]
Using Taylor expansion, we deduce that
$\frac{a^\prime}{a}(s)\geq (1-\delta)\frac{b^\prime}{b}(s)$ for all $s>T_0$ 
which are sufficiently close to $T_0$. This  
contradicts the choice of $T_0$. 
Since $\delta >0$ was arbitrary, the lemma follows.
\end{proof}





A \emph{cusp} $T$ is an unbounded component of $M_{\rm thin}$. 
It is homeomorphic to $N\times [0,\infty)$ where
$N$ is a closed manifold of dimension $n-1$. 
The manifold $N$ is 
homeomorphic to the quotient of 
a horosphere in the universal covering $\tilde M$ of $M$ 
by a parabolic subgroup of the isometry group of 
$\tilde M$. 
As before, the boundary of $T$ need not be smooth, but
everything said so far for boundaries of Margulis tubes
is also valid for cusps (see the remark after Theorem 2.14 in 
\cite{BCD93}). In particular, 
Theorem \ref{thmbcd93} holds true for boundaries of cusps.





A version of Lemma \ref{volumegrowth} is also valid for cusps and follows with exactly the same
argument. Namely, let $T$ be a cusp, with boundary $\partial T$. Write $T=\partial T\times [0,\infty)$ where
$\partial T\times \{s\}$ is the hypersurface of distance $s$ to $\partial T$. 
Let $\eta:[0,\infty)\to T$ be a radial
geodesic and let $J(t)=(J_1(t),\dots,J_{n-1}(t))$ be Jacobi fields with the following properties. 
The vectors $J_1(0),\dots,J_{n-1}(0)$ define an orthonormal basis of 
$\eta^\prime(0)^\perp$, and $\Vert J(t)\Vert \to 0$ $(t\to \infty)$. Denote by $j(t)$ the 
determinant of $J(t)$, viewed as a matrix with respect to 
an orthonormal basis of $\eta^\prime(t)^\perp$ defined by parallel vector fields along $\eta$.

 Let $\bar j(t)$ be the corresponding function for Jacobi fields defined by a horoball in hyperbolic 
 $n$-space. Explicit calcuation yields $\bar j(t)=e^{-(n-1)t}$ and hence
 $\bar j(T)= (n-1)\int_T^\infty \bar j(s)ds$ for all $T$. 
 The comparison argument from the proof of 
 Lemma \ref{volumegrowth} together with the results in \cite{E87}  
 imply that the inequality $j(T)\geq \int_T^\infty j(s)ds$ holds true for cusps in 
 a manifold $M$ of curvature $\leq -1$. In fact, this statement can also be obtained as a limiting
 case  of Lemma \ref{volumegrowth} by reparametrizing 
the radial geodesic arc $\eta$ of the Margulis tube, 
choosing $\eta(R)$ as a basepoint, renormalizing 
the Jacobi fields with rescaling and the Gram-Schmidt procedure 
and letting $R$ tend to infinity. 


As one consequence of the above discussion, 
by possibly replacing $M_{\rm thick}$ by 
a neighborhood of uniformly bounded radius
we may assume that 
\begin{equation}\label{volume}
\frac{3}{2} {\rm vol}(M_{\rm thick})\geq  {\rm vol}(M).\end{equation}






Consider again a Margulis tube $T$ in $M$.
Recall that the boundary
$\partial T$ of $T$ equals the set $\{\rho=R\}$ which can 
be viewed as a graph over the unit normal bundle $N(\gamma)$ of $\gamma$
for some smooth function $R:N(\gamma)\to (0,\infty)$.
Here we use Theorem \ref{thmbcd93} to assure that the function $R$ is smooth. 
We equip $\partial T$ with the volume element determined by this description
(this is in general not the volume element of $\partial T$ viewed as
a smooth hypersurface in $M$). By this we mean that the volume 
element of the hypersurface $\partial T$ at a point 
$(\sigma_0,s_0,R(\sigma_0,s_0))$ equals the radial 
projection of the volume 
element of the local hypersurface $\{(\sigma,s,R(\sigma_0,s_0))\}$ 
passing through $(\sigma_0,s_0,R(\sigma_0,s_0))$.  
This makes sense since by  Theorem \ref{thmbcd93}, 
radial geodesics intersect $\partial T$ transversely. 
Or, equivalently, 
this volume element is chosen in such a way that the Jacobian at 
$(\sigma,s)$ of the map 
$(\sigma,s)\to (\sigma,s,R(s,\sigma))$  equals
the Jacobian of the normal exponential map $(\sigma,s)\to 
\exp(R(\sigma,s)(\sigma,s))$, and this is just the function $j$ 
from Lemma \ref{volumegrowth}. 
The same construction is also valid for a cusp $T$, 
and we equip $\partial T$ with the corresponding volume element
obtained by projection of the volume element on local 
hypersurfaces orthogonal to the radial geodesics. 
 

\begin{lemma}\label{partialintegration}
Let $T\subset M_{\rm thin}$ be a Margulis tube or a cusp with boundary 
$\partial T$, and 
let $f$ be a smooth function
on $T$ with $\int_T f^2\geq \int_{\partial T}f^2$; then 
\[\int_T f^2\leq \frac{4}{(n-2)^2}\int_T \Vert \nabla f \Vert^2.\]
\end{lemma}
\begin{proof} We begin with the case that $T$ is a Margulis tube. 
Let $\gamma:[0,\ell)\to M$ be a parameterization of the core geodesic 
of $T$ by arc length. We use normal exponential coordinates and Lemma \ref{volumegrowth}.
Let $\rho$ be the radial distance from $\gamma$ and let 
$j(\sigma,s,\rho)$ be the Jacobian of the normal exponential map 
at the point with
"coordinates" $(\sigma,s,\rho)$. 
Then we have 
\[\int_Tf^2=\int_{S^{n-2}}d\sigma\int_0^\ell ds
\int_0^{R(\sigma,s)} f^2 j(\sigma,s,\rho) d\rho.\]


Define $a(\sigma,s,\rho)=\int_0^{\rho}j(\sigma,s,u)du$. 
By Lemma \ref{volumegrowth} 
and the definition of the volume element on $\partial T$, 
integration by parts along the radial rays from $\gamma$ 
yields 
\begin{equation}\label{inequality1}
\int_T f^2\leq \frac{1}{n-1}
\int_{\partial T} f^2 - 
2\int_{S^{n-2}}d\sigma\int_0^\ell ds \int_0^{R(\sigma,s)} ff^\prime 
a(\sigma,s,\rho) d\rho\notag \end{equation}
where $f^\prime$ is the derivative of $f$ in direction of the 
radial variable $\rho$.

 
By the assumption in the lemma, 
we have 
\[\int_{\partial T} f^2 \leq \int_Tf^2\]
and therefore taking absolute values 
and using Lemma \ref{volumegrowth} once more, 
we obtain
\begin{align}
\frac{n-2}{n-1}\int_T f^2  & \leq  
\frac{2}{n-1}
\int_{S^{n-2}}d\sigma\int_0^\ell ds \int_0^R \vert  ff^\prime \vert j(\sigma,s,\rho) d\rho
\notag\\ & =
\frac{2}{n-1}\int_T \vert  ff^\prime \vert 
\leq \frac{2}{n-1}
(\int_T f^2)^{1/2}(\int_T \Vert \nabla f\Vert^2)^{1/2} \notag\end{align}
where the last inequality follows from 
Schwarz's inequality and from $\vert f^\prime\vert \leq 
\Vert \nabla f\Vert$ (compare the proof of  
Lemma 1 of \cite{DR86}). Dividing by 
$\frac{2}{n-1}(\int_T f^2)^{1/2}$ and 
squaring the resulting inequality yields the lemma in the case that
$T$ is a Margulis tube.

The argument for a cusp is almost identical.
Thus let $T$ be such a cusp, with boundary $\partial T$. 
Then $\partial T$ is an closed manifold, 
and $T$ is diffeomorphic to $\partial T\times [0,\infty)$.
Let $\rho:T\to [0,\infty)$ be the radial distance from $\partial T$. 
 

Let 
$f:T\to \mathbb{R}$ be a smooth square integrable function with 
$\int_Tf^2\geq \int_{\partial T}f^2$ and let 
$j:\partial T\times [0,\infty)\to 
(0,\infty)$ be the function which describes the volume form on $T$
as discussed above. The version of 
Lemma 2.1 for cusps and integration by parts is used as in the proof for
Margulis tubes.
Since $f$ is square integrable,
as $R\to \infty$ we have $\int_{\rho=R}f^2\to 0$ and hence
the same calculation as before yields 
(here $d\mu$ is the 
volume element on $\partial T$
 used in the lemma) 
\begin{align} \int_Tf^2 &=\int_{\partial T}d\mu
\int_{0}^\infty f^2 j(x,\rho)d\rho
=\lim_{R\to\infty}
\int_{\partial T}d\mu\int_{0}^R f^2j(x,\rho)d(\rho)\notag \\
&\leq \frac{1}{n-1}\bigl(
\int_{\partial T}f^2+2\int_{\partial T}d\mu
\int_0^\infty ff^\prime j(x,\rho) d\rho \bigr).\notag
\end{align}  
As before,   
this yields the required estimate. 
Note that the change of sign in this formula stems from the fact that
here $\rho$ is the radial distance from the boundary, while for Margulis
tubes, $\rho$ denoted the radial distance from the core geodesic which 
seems more natural in that context. 
\end{proof}


 

For a Margulis tube $T$ or a cusp, let $\hat T$ be the set of all points 
$x \in T$ whose radial distance from the boundary $\partial T$ is at least one. 
Thus if $T$ is Margulis tube, determined by a closed
geodesic and a function $R(\sigma,s)>0$ on the normal bundle of that geodesic, then 
using "coordinates" $(\sigma,s,\rho)$ as before,  
we have  $\hat T=\{(\sigma,s,\rho)\in T\mid \rho\leq R(\sigma,s)-1\}$. 
By our assumption on Margulis tubes,
$T-\hat T$ is diffeomorphic to $\partial T\times [0,1)$.
Write $\hat M_{\rm thick}=M-\cup \hat T$ where the union is over all 
Margulis tubes and cusps. 


For a smooth square integrable function $f$ on $M$ denote by
\[{\cal R}(f)=\int_M \Vert \nabla f\Vert^2/\int_M f^2\]
the Rayleigh quotient of $f$. 

 
\begin{lemma}\label{rayleigh}
Let $f:M\to \mathbb{R}$
be a smooth square integrable function with
Rayleigh quotient ${\cal R}(f)<  \frac{(n-2)^2}{12}$; then 
\[\int_{\hat M_{\rm thick}}f^2\geq \frac{1}{3}\int_M f^2.\]
\end{lemma}
\begin{proof}
By our assumption on the components of the thin part of $M$, the set
\[A=\hat M_{\rm thick}-M_{\rm thick}\]  is a union of \emph{shells},
i.e. submanifolds of $M$ with boundary which either are
diffeomorphic to $S^{n-2}\times S^1\times [0,1]$
(if the component is a Margulis tube) or to 
$N\times [0,1]$ (if the component is a cusp),
where
$N$ is a quotient of a horosphere by a rank
$n-1$ parabolic subgroup of the isometry group of 
the universal covering $\tilde M$ of $M$. 


Let $f:M\to \mathbb{R}$ be a smooth function with 
${\cal R}(f)<\frac{(n-2)^2}{12}$. We want to show that
\begin{equation}\label{integral}
\int_{\hat M_{\rm thick}}f^2\geq 
\frac{1}{3}\int_M f^2,\notag\end{equation} 
and to this end we assume to the contrary 
that $\int_{\hat M_{\rm thick}}f^2<
\frac{1}{3}\int_M f^2$.

Recall that $M_{\rm thin}$ is a disjoint union 
of a finite number of 
Margulis tubes and cusps, say 
$M_{\rm thin}=\cup_{i=1}^kT_i$.
For each of these tubes and cusps $T_i$, let $r_i$ be the radial distance
function to the boundary hypersurface, i.e. $r_i(x)$ is 
the length of the radial arc
connecting the point $x\in T$ to $\partial T$. By reordering we may assume that 
there exists a number $p\leq k$ such that 
for all $i\leq p$, there is some 
$s_i\leq 1$ such that
\[\int_{\{r_i=s_i\}\cap T_i}f^2\leq \int_{T_i\cap \{r_i\geq s_i\}}f^2\]
and that for $i>p$, such an $s_i$ does not exist.
Here we use the volume element on the hypersurfaces
$\{r_i=s_i\}$ as in Lemma \ref{partialintegration}. 

We distinguish two cases. In the first case,
$\sum_{i=1}^p\int_{T_i-A}f^2\geq \frac{1}{3}\int_Mf^2$.
Lemma \ref{partialintegration} then shows that
\begin{align}
\int_M\Vert \nabla f\Vert^2\geq 
\sum_{i=1}^p\int_{T_i\cap \{r_i\geq s_i\}}\Vert \nabla f\Vert^2 \notag \\
\geq \frac{(n-2)^2}{4}
\sum_{i=1}^p\int_{T_i\cap \{r_i\geq s_i\}}f^2\geq 
\frac{(n-2)^2}{12}\int_M f^2.\notag\end{align}
Thus we have
${\mathcal R}(f)\geq \frac{(n-2)^2}{12}$ 
which contradicts our assumption on $f$.

In the second case, we have
$\sum_{i=1}^p\int_{T_i-A}f^2< \frac{1}{3}\int_Mf^2$.
Then 
\begin{equation}\label{eqlemma42}
\sum_{i=p+1}^k\int_{T_i-A}f^2\geq \frac{1}{3}\int_M f^2.
\end{equation} 
But for each $i>p$, integration of the defining equation
\begin{equation}\label{integral2}
\int_{T_i\cap \{r_i=s\}}f^2\geq 
\int_{T_i\cap \{r_i\geq s\}}f^2\end{equation}
over the shell $T_i\cap A=\{0\leq r_i\leq 1\}$ yields
\[\int_0^1ds \int_{T_i\cap \{r_i=s\}}f^2=
 \int_{T_i\cap A}f^2\geq \int_{T_i-A}f^2.\]
Summing over $i\geq p+1$ and using inequality
(\ref{eqlemma42}), we obtain
\[\int_{\cup_{i=p+1}^kT_i\cap A}f^2\geq \sum_{i=p+1}^k\int_{T_i-A}f^2
\geq \frac{1}{3}\int f^2.\]
As $A\subset \hat M_{\rm thick}$, this 
contradicts the assumption on $f$. The
lemma follows. 
\end{proof}

The next proposition completes the proof of Theorem \ref{thm1} for a choice of a Margulis
constant so that $\hat M_{\rm thick}$ as defined above is contained in the
thick part of $M$ for this constant. 
   

\begin{proposition}\label{lowerbound}
For all $k\geq 0$ we have 
\[\lambda_k(M)\geq \min\{
\frac{1}{3}\lambda_k(\hat M_{\rm thick}),
\frac{(n-2)^2}{12}\}\]
for any finite volume oriented Riemannian mani\-fold $M$ of dimension
$n\geq 3$ and curvature $\kappa\in [-b^2 -1]$.  
\end{proposition}
\begin{proof} Let $M$ be a finite volume oriented Riemannian manifold of dimension
$n\geq 3$ and curvature $\kappa \leq -1$.
We use the previous notations.

If $M$ is non-compact, then by Corollary 3.2 of \cite{H04}, 
the bottom of the essential 
spectrum of $M$ is not smaller than $(n-1)^2/4$. Thus the intersection of the spectrum of 
$M$ with the interval $[0,\frac{(n-2)^2}{12}]$ consists of a finite number of eigenvalues. 


Let ${\cal H}(M)$ be the Sobolev space 
of square integrable functions on $M$, with square integrable weak 
derivatives. 
Let $k>0$ be such that $\lambda_{k}(M)<(n-2)^2/12$. 
Let $m\leq k-1$ be the largest number so that $\lambda_m(M)<\lambda_k(M)$.
Note that as we count eigenvalues with multiplicities,
we may have $m<k-1$. Choose a $k-m$-dimensional subspace 
$E_k\subset {\cal H}(M)$ of the eigenspace for the eigenvalue $\lambda_k(M)$ and define 
$E=V\oplus E_k$ where $V$ is the direct sum
of the eigenspaces for eigenvalues strictly smaller than $\lambda_k(M)$.
In particular, 
$E$ is a $k+1$-dimensional linear subspace of the Hilbert space ${\cal H}(M)$. 


By construction, $\hat M_{\rm thick}$ is a smooth manifold with smooth boundary. 
Denote by ${\cal H}(\hat M_{\rm thick})$ the Sobolev space 
of square integrable functions
on $\hat M_{\rm thick}$ with square integrable weak derivatives. 
Here the weak derivative of a function $f$ on $\hat M_{\rm thick}$
is a vector field $Y$ so that
\[\int_{\hat M_{\rm thick}} \langle Y,X\rangle=
-\int_{\hat M_{\rm thick}}f{\rm div}(X)\]
for all smooth vector fields $X$ on $\hat M_{\rm thick}$ with compact
support in the interior of $\hat M_{\rm thick}$.
%If we denote by $\nu$ the outer normal field of 
%$\partial \hat M_{\rm thick}$, then 
Green's formula implies that
${\cal H}(\hat M_{\rm thick})$ contains all functions 
$f$ on $\hat M_{\rm thick}$ which are smooth up to and including the 
boundary.


As smooth functions are dense in ${\cal H}(M)$ and 
${\cal H}(\hat M_{\rm thick})$, there 
is a natural linear one-Lipschitz restriction map
$\Pi:{\cal H}(M)\to {\cal H}(\hat M_{\rm thick})$. We denote by 
$W$ the image of the linear subspace $E$ under $\Pi$. 

We claim that
${\rm dim}(W)=k+1$. To this end assume otherwise. Then there is 
a normalized 
function $f\in E$ (i.e. $\int_M f^2=1$) 
whose restriction to $\hat M_{\rm thick}$ vanishes.
But by the definition of $E$, the Rayleigh quotient 
${\cal R}(f)$ of $f$ is at most $\lambda_k(M)<(n-2)^2/12 $ and therefore 
Lemma \ref{rayleigh} shows that $\int_{\hat M_{\rm thick}}f^2\geq \frac{1}{3}\int_M f^2$.

As ${\rm dim}(W)=k+1$ and as 
${\cal H}(\hat M_{\rm thick})$ is the function space for 
$\hat M_{\rm thick}$ with Neumann boundary conditions (see p.14-17 in \cite{C84}),  
Rayleigh`s theorem shows that 
there exists a normalized function $f\in E$ 
with ${\cal R}(\Pi(f))={\cal R}(f\vert \hat M_{\rm thick})
\geq\lambda_k(\hat M_{\rm thick})$.
Furthermore, as $f\in E$, we have 
${\cal R}(f)\leq \lambda_k(M)$. Note that $f$ and $\Pi(f)$ are smooth. 

By Lemma \ref{rayleigh}, we have
$\int_{\hat M_{\rm thick}}f^2\geq \frac{1}{3}\int_M f^2=\frac{1}{3}$ and hence
\[\lambda_k(M)\geq {\cal R}(f)\geq \int_{\hat M_{\rm thick}}\Vert \nabla f\Vert^2 
\geq \frac{1}{3}{\cal R}(f\vert_{\hat M_{\rm thick}})\geq 
\frac{1}{3}\lambda_k(\hat M_{\rm thick}).\] 
This is what we wanted to show.
\end{proof}


As an easy consequence, we obtain the estimate of Schoen \cite{S82}.

\begin{corollary}\label{eigenbelow}
For every $n\geq 3$ and every $b\geq 1$ 
there exists a number $\chi(n,b)>0$ such that 
\[\lambda_1(M)\geq \frac{\chi(n,b)}{{\rm vol}(M)^2}\]
for every finite volume Riemannian manifold $M$ of 
dimension $n$ and curvature
$\kappa\in [-b^2,-1]$. 
\end{corollary}
\begin{proof} Let $\epsilon >0$ be a Margulis constant for Riemannian 
manifolds $M$ of dimension $n\geq 3$ and curvature in the interval
$[-b^2,-1]$. As $M_{\rm thick}\not=\emptyset$, the manifold $M$ contains
an embedded ball of radius $\epsilon$ which is isometric to a ball of the
same radius in a simply connected manifold of curvature contained in 
$[-b^2,-1]$. By comparison, the volume of such a ball is bounded from 
below by a universal constant $a(n,b)$ only depending on $n$ and $b$ and hence
the volume of every Riemannian $n$-manifold of curvature in $[-b^2,-1]$ is bounded 
from below by 
$a(n,b)$. 
Thus Theorem \ref{thm1} shows that 
$\lambda_1(M)\geq \min\{a(n,b)^2(n-2)^2/12{\rm vol}(M)^2,
\lambda_1(M_{\rm thick})/3\}$.


The manifold with boundary $M_{\rm thick}$ 
is uniformly quasi-isometric to a finite connected graph $G$,
and its first nontrivial eigenvalue 
$\lambda_1(M_{\rm thick})$ with Neumann boundary conditions
is uniformly equivalent to the first nontrivial eigenvalue of the 
graph Laplacian \cite{Man05} (see also Lemma 2.1 of \cite{LS12} for an explicit
formulation of this fact). 


Now the first eigenvalue $\lambda_1(G)$ 
of a finite graph $G$ satisfies $\lambda_1(G)\geq h^2(G)/2$ where
$h(G)$ is the so-called \emph{Cheeger constant} of the graph \cite{Chu96}.
This Cheeger constant is defined to be 
\[\min\frac{\vert E(U,U^\prime)\vert}{\min\{\vert U\vert,\vert U^\prime\vert\}}\]
where $U$ is a subset of the vertex set ${\mathcal V}(G)$ of $G$, 
$U^\prime={\mathcal V}(G)-U$ and where $E(U,U^\prime)$ is the set of 
all edges which connect a vertex in $U$ to a vertex in $U^\prime$. 
As $G$ is connected, 
this Cheeger constant is at least $1/\vert {\cal V}(G)\vert \sim
1/{\rm vol}(M)$ which yields the corollary. 
\end{proof}





\section{Bounding small eigenvalues from above}\label{above}

In this section we restrict to the investigation of orien\-ted 
hy\-per\-bolic 
3-manifolds of finite volume. 
Our goal is to prove Theorem \ref{thm2}. 


 
We begin with analyzing in more detail functions on Margulis tubes 
in such a hyperbolic $3$-manifold $M$. 
For a sufficiently
small Margulis constant
$\epsilon>0$, the boundary $\partial T$ of each such tube $T$ is a flat
torus whose injectivity radius for the induced metric 
roughly equals $\epsilon$ \cite{BCD93}. 
Furthermore, radial geodesics intersect $\partial T$ orthogonally. 
Note that an analogous statement is
not true in higher dimensions.

The \emph{radius} ${\rm rad}(T)$ 
of the tube $T$ is the distance of $\partial T$ to 
the core geodesic $\gamma$. The following lemma is only valid
in dimension three. 


\begin{lemma}\label{radiusbound}
There is a number $q_1=q_1(\epsilon)>0$ only depending on $\epsilon$ such that
${\rm rad}(T)\geq \log {\rm vol}(\partial T)-q_1$. 
\end{lemma}
\begin{proof} 
Write $R={\rm rad}(T)$. We may assume that $R>1$.
The tube $T$ is isometric to the quotient of the tubular
neighborhood $N(\tilde \gamma,R)$ of radius $R$ 
of a geodesic $\tilde \gamma$ in $\mathbb{H}^3$ under
an infinite cyclic group of loxodromic isometries. Up to conjugation, the generator
$\psi$ of this group 
is determined by its complex translation length $\chi\in \mathbb{C}$
with $\Re(\chi)=\ell<2\epsilon$. Here $\ell$ is the length of the core geodesic of $T$.


The boundary $\partial N(\tilde \gamma, R)$ of $N(\tilde \gamma,R)$ 
is a flat two-sided infinite cylinder of circumference $2\pi \sinh R$. 
For a fixed identification of the fibre
of the unit normal circle bundle of $\tilde \gamma$ over the point
$\tilde \gamma(0)$ with the unit circle $S^1$, parallel transport along $\tilde\gamma$ and the 
normal exponential map determine global ``coordinates'' on 
$\partial N(\tilde \gamma,R)$. These ``coordinates'' consist in a 
diffeomorphism 
$S^1\times \mathbb{R}\to \partial N(\tilde \gamma,R)$. The isometry $\psi$   
identifies the meridian
$\exp(R S^1\times \{0\})$ 
with the meridian $\exp(R S^1\times \{\ell\})$   
by an isometry
which is given by rotation with angle $\Im \chi$
in these coordinates. In particular, 
the volume of the boundary torus $\partial T$ of $T$ 
equals 
\begin{equation}\label{volumeformula}
{\rm vol}(\partial T)=2\pi\ell \sinh R\cosh R
\end{equation} 
and is independent of 
$\Im \chi$. 

The second fundamental form of the torus 
$\partial T$ is uniformly bounded, independent of $R>1$. 
This implies  that 
the length of
the shortest geodesic on $\partial T$ is contained
in an interval $[\epsilon,a\epsilon]$ where $a>0$ is a
universal constant (compare Theorem 2.14  and the 
following discussion in \cite{BCD93}).
As we are only interested in tubes of large radius, 
we may assume that the length
$2\pi \sinh R$ of the meridian of 
$\partial T$ is bigger than $a\epsilon$. 

Now the distance on $\partial N(\tilde \gamma,R)$ 
for the intrinsic path metric between the 
circles corresponding to the coordinates $s=0,s=\ell$ 
equals $\ell\cosh R$. Here as before, $s$ is the length parameter
of $\gamma$. A shortest closed geodesic
on the flat torus 
$\partial T$ lifts to a straight line segment on
$\partial N(\tilde \gamma,R)$. Since $2\pi\sinh R>a\epsilon$, 
this line segment connects the circle $\{s=0\}$ to the circle
$\{s=k \ell\}$ for some integer $k\geq 1$. From this we conclude that
\[\ell\cosh R\leq a\epsilon.\]

From the formula (\ref{volumeformula})
for ${\rm vol}(\partial T)$, we deduce that
$\sinh R\geq \frac{1}{2\pi a\epsilon}{\rm vol}(\partial T)$ and hence 
\[e^R\geq \frac{1}{\pi a\epsilon}{\rm vol}(\partial T)\]
which is what we wanted to show.
\end{proof}

\begin{remark} 
For hyperbolic manifolds of dimension $n\geq 4$, 
Proposition 2 of \cite{BS87} states a reverse inequality:
Namely, if $U$ is any Margulis tube, and if $r(U)$ is
the largest distance of a point in $U$ to the boundary,
then ${\rm Vol}(U)\geq d_n\sinh(\frac{1}{3}r(U))$
where $d_n>0$ is a constant only depending on $n$.
The case ${\rm dim}(M)=3$ is very special as every point
on the boundary of a Margulis tube has the same distance
to the core geodesic. We refer to p.3 of \cite{BCD93} for more
information.
\end{remark} 

In the statement of the next proposition, we use the fact that 
given a Margulis tube $T$ of radius at least three, the volume of 
$T$ is uniformly proportional to the volume of its boundary 
$\partial T$. 


\begin{proposition}\label{tube}
There exists a number $q_2=q_2(\epsilon)>0$ with the following
property. Let $T\subset M$ be a Margulis tube or a cusp 
with boundary $\partial T$.
Let $f:\partial T\to \mathbb{R}$ be a function 
(not neccessarily of zero mean) whose
Rayleigh quotient equals $d\geq 0$. 
Then there is an extension of $f$ to a smooth function $F$ on $T$ with 
the following properties.
\begin{enumerate}
\item $\frac{1}{4}\int_{\partial T}f^2\leq 
\int_TF^2\leq \frac{1}{2}\int_{\partial T}f^2.$
\item The 
Rayleigh quotient of $F$ is at most
$dq_2\log {\rm vol}(T)$. 
\item If $\int_{\partial T}f=0$ then $\int_T F=0$.
\end{enumerate}
\end{proposition}
\begin{proof} Let $\gamma$ be the core curve of the tube $T$.
Use the coordinates 
$(\sigma,s,\rho)$ on $T$, given by the angular coordinate
$\sigma$ on the unit normal circle over a point in $\gamma$, the length
parameter $s$ on $\gamma$ and the distance $\rho$ from $\gamma$. 
If $R>0$ is the radius of $T$ then 
the boundary $\partial T$ of $T$ is the surface
$\rho=R$. This boundary is a flat torus whose injectivity
radius equals $\epsilon$ up to a universal multiplicative
constant. 

%Let $k>0$ be the volume of $\partial T$. 
By Lemma \ref{radiusbound}, 
the radius $R$ of the tube satisfies
\[-\theta=\log{\rm vol}(\partial T)  -q_1-1-R\leq -1\] where 
$q_1>0$ is a universal constant.  By the inequality
(\ref{mincontrol}) in Section 2, we 
may assume that 
$R-\theta \geq 3$.

Let $f:\partial T=\{\rho=R\}\to \mathbb{R}$ be a smooth function
with Rayleigh quotient $d\geq 0$. 
Decompose $f=f_0+g$ 
where $f_0$ is a constant function and $\int_{\partial T} g=0$. 
Extend $f_0$ to a constant
function $F_0$ on $T$, and extend the function $g$  
to a function $G:T\to \mathbb{R}$ by
\begin{equation}
G(\sigma,s,\rho)=
\begin{cases}
g(\sigma,s,R) &\text{if $\theta+1\leq \rho\leq R$;}\\
(\rho-\theta) g(\sigma,s,R) &\text{if 
$\theta\leq \rho\leq \theta +1$;}\\
0 &\text{otherwise.}
\end{cases}
\end{equation}
Then $F=F_0+G$ is continuous, smooth away from the hypersurfaces
$\rho=\theta+1$ and $\rho=\theta$, with uniformly bounded derivative.
Standard Sobolev theory then implies that $F$ is 
contained in the Sobolev space
of square integrable functions with square integrable weak derivative,
and $F\vert \partial T=f$.

As for each $r\in (0, R)$ 
the radial projection $(\sigma,s,r)\to (\sigma,s,R)$ 
of the torus $\{\rho =r\}$ onto the boundary torus $\partial T$ 
is a homothety, with dilation
$\sinh(R)\cosh(R)/\sinh(r)\cosh(r)$, we have $\int_TG =0$. 
Furthermore, 
integration using the explicit description of the metric on $T$
as a warped product metric (see Section \ref{below} for details) yields that
\begin{equation}\label{squarenorm}
\int_T (F_0+G)^2\in [\frac{1}{4}\int_{\partial T} (f+g)^2,\frac{1}{2}\int_{\partial T}(f+g)^2].
\end{equation}
Namely, as $(F_0+G)^2(\sigma,s,\rho)\leq f^2(\sigma,s,R)$ for all $\rho$, 
the upper bound follows from
\[\int_T(F_0+G)^2\leq \int_{\partial T}f^2\int_\theta^R\frac{\sinh \rho \cosh \rho}{\sinh R\cosh R}d\rho
=\frac{\sinh^2(R)-\sinh^2(\theta)}{2\sinh R\cosh R}\int_{\partial T}f^2.\]
To establish the lower bound, simply note that
\begin{align}
\int_T(F_0+G)^2\geq \int_{\partial T}f^2\int_{\theta+1}^R 
\frac{\sinh \rho \cosh \rho}{\sinh R\cosh R}d\rho\notag\\
=\frac{\sinh^2(R)-\sinh^2 (\theta+1)}{2\sinh R\cosh R}\int_{\partial T}f^2.
\notag\end{align}
Up to adjusting the requirement on the radius of the Margulis tube, we may
assume that $\sinh^2(\theta +1)\leq \frac{1}{4} \sinh(R)^2$ and
$\frac{3}{4}\sinh(R)\geq \frac{1}{2}\cosh(R)$
which then results in the 
estimate stated in the first part of the proposition. 

Now $\int_T\Vert \nabla(F_0+G)\Vert^2=\int_T\Vert \nabla G\Vert^2$
and therefore for the second part of the statement in the proposition, 
it suffices to show that
\begin{equation}\label{logestim}
  \int_T\Vert \nabla G\Vert^2\leq c_1(R-\theta) 
  \int_{\partial T}\Vert \nabla g\Vert^2\end{equation}
for a universal 
constant $c_1>0$. Namely, the discussion in the previous paragraph shows
that the integral of the constant function one over the shell
$\theta\leq \rho\leq R$ is bounded from below by 
$\frac{1}{4}{\rm vol}(\partial T)$. In particular, we have 
$\log{\rm vol}(T)\geq \log {\rm vol}(\partial T)-\log 4$. 
On the other hand, by the choice of 
$\theta$, we also have $R-\theta=\log{\rm vol}(\partial T)-q_1-1$. 

To establish the estimate (\ref{logestim})
recall from Chapter 2 of \cite{C84} that the first 
nonzero eigenvalue of $\partial T$ is not smaller than
$c_2/{\rm vol}(\partial T)^2$ where $c_2>0$ is a universal constant
(here we use that the injectivity radius of $\partial T$ is at least $\epsilon$).
In particular, by the definition of the number $\theta$, this eigenvalue is 
not smaller than $c_3e^{-2(R-\theta)}$ where $c_3>0$ is a universal constant.
As a consequence, we have 
\begin{equation}\label{rayleighbd}
\int_{\partial T}\Vert \nabla g\Vert^2\geq c_3e^{-2(R-\theta)}\int_{\partial T}g^2.
\end{equation}

Now for $r\in [\theta+1,R]$, the radial projection of 
the torus $\{\rho=r\}$ onto $\partial T$ 
scales the metric with a fixed constant. Moreover,
the directional derivative of $G$ in direction of the
radial vector field vanishes. This implies that   
\[\int_{\rho=r}\Vert \nabla G(\sigma,s,r)\Vert^2 =
\int_{\partial T} \Vert \nabla G(\sigma,s,R)\Vert^2\]
for $r\in [\theta +1,R]$ and hence
\[\int_{\theta+1\leq \rho\leq R} \Vert \nabla G(s,t,\rho)\Vert^2=
(R-\theta-1)\int_{\partial T}\Vert \nabla g\Vert^2.\]
On the other  hand, if
$\theta< r<\theta+1$ then 
\[\int_{\rho=r}\Vert \nabla G(s,t,\rho)\Vert^2 \leq 
\int_{\partial T}\Vert \nabla G(s,t,R)\Vert^2+\int_{\rho=r}G(s,t,\rho)^2.\]
Here the second term in this inequality is the contribution of the
derivative of $G$ in direction of the radial vector field, and we use
that this vector field is normal to the hypersurfaces
$\{\rho=r\}$. 

There is a universal constant $c_4>0$ such that for each $r\in 
[\theta,\theta+1]$ the volume of the level surface
$\{\rho=r\}$  is not bigger than 
$c_4e^{-2(R-\theta)}({\rm vol}(\partial T))$. Integration of this
inequality over the interval $[\theta,\theta +1]$ yields
$\int_{\theta\leq \rho\leq \theta +1}G^2\leq
c_4e^{-2(R-\theta)}\int_{\partial T}g^2$. Hence by the estimate
(\ref{rayleighbd}), 
\[\int_{\theta\leq \rho\leq \theta+1}G^2\leq 
\frac{c_4}{c_3}\int_{\partial T}\Vert \nabla g\Vert^2.\]
Together this
implies
\begin{equation}\label{secondbound}
\int_T\Vert \nabla G\Vert^2\leq (R-\theta +\frac{c_4}{c_3})
\int_{\partial T} \Vert \nabla g\Vert^2. \end{equation}
The estimates (\ref{squarenorm},\ref{secondbound}) 
together with the choice of $\theta$ show 
part (1) and (2) of the proposition. The third part is immediate from the 
fact that by construction, if $\int _{\partial T}f=0$ then $F=G$ and 
$\int_{T}G=0$.

The case that $T$ is a cusp is completely analogous but easier 
and will be omitted.
\end{proof} 



As an immediate consequence of Proposition \ref{tube} we obtain

\begin{corollary}\label{neumanntube}
The first nonzero eigenvalue with Neumann boundary conditions
of a three-dimensional hyperbolic Margulis tube or a cusp is at most
$q_3\log{\rm vol}(T)/{\rm vol}(\partial T)^2$
where $q_3>0$ is a universal constant. 
\end{corollary}
\begin{proof}
By the discussion in 
Chapter II, Section 2 of \cite{C84} 
and the fact that the injectivity radius of the flat torus
$\partial T$ is proportional to $\epsilon$,  
there exists a universal number $a>0$ and 
a function $f:\partial T\to \mathbb{R}$ with $\int_{\partial T}f=0$ and 
\[\int_{\partial T}\Vert \nabla f\Vert^2\leq a\int_{\partial T}f^2 /{\rm vol}(\partial T)^2.\]

By Proposition \ref{tube}, the function $f$ can be extended to a function
$F:T\to \mathbb{R}$ with $\int_T F=0$ and Rayleigh quotient
${\cal R}(f)\leq q_2a\log {\rm vol}(T)/{\rm vol}(\partial T)^2$. 
Now $F$ is smooth up to and including the boundary, and therefore 
by Rayleigh's principle as 
explained in Chapter 1 Section 5 of \cite{C84},
this implies 
that the first non-zero eigenvalue of $T$ with Neumann boundary conditions 
is bounded from above by $q_2a\log{\rm vol}(T)/{\rm vol}(\partial T)^2$ as
claimed in the corollary. 
\end{proof}


%\begin{remark} We do not know whether the estimate in 
%Corollary \ref{neumanntube} is sharp.
%\end{remark}

The strategy for the proof of 
Theorem \ref{thm2} from the introduction consists in extending
an eigenfunction $f$ on $\hat M_{\rm thick}$ to a function on $M$
using the construction in Proposition \ref{tube}. But the
control on square integrals and Rayleigh quotients
established in Proposition \ref{tube} depends on a control of
these data for the restriction of $f$ to the boundary of
$\hat M_{\rm thick}$, and there is no apparent relation to the
global square norm and the global Rayleigh quotient of $f$.
Note also that there is no useful Harnack inequality for the restriction of 
an eigenfunction on $\hat M_{\rm thick}$ to its boundary. 

To overcome this difficulty we establish some a-priori control
on the restriction of eigenfunctions for small eigenvalues to sufficiently large
tubular neighborhoods of the boundary of $\hat M_{\rm thick}$- the price
we have to pay is that we have to decrease the Margulis constant
which appears in Theorem \ref{thm2}. 

We begin with establishing 
a modified version of Lemma \ref{partialintegration}.
We only formulate this lemma for hyperbolic 3-manifolds although it 
is valid for arbitrary finite volume manifolds of curvature contained in 
$[-b^2,-1]$, where integration over radial distance hypersurfaces has to be
interpreted as in Section \ref{below}. 

Define a \emph{shell}
in a Margulis tube or cusp $T$ to be a subset of $T$ of the form
$\{s\leq \rho\leq t\}$ where $\rho$ is up to an additive  
constant the \emph{negative} of the
radial distance from the boundary of $T$
and $s<t$ (using the negative of the radial distance is for convenience here). 
The \emph{height} of the shell equals $t-s$. 
We always write a shell $N$ in the form $N=V\times [0,k]$ where
$k$ is the height of the shell, the manifold
$V$ is diffeomorphic to the boundary 
$\partial T$ of the tube or cusp, 
the real parameter $t$ is the radial distance
from the boundary component $V\times \{0\}$, and 
the boundary component $V\times \{k\}$ is closer to the boundary $\partial T$ of the tube or cusp. 
As before, we denote by ${\cal R}(f)$ the Rayleigh quotient of a function $f$.


\begin{lemma}\label{volumecontrol}
\begin{enumerate}
\item Let $k\geq 2$ and let 
$N=V\times [0,k]$ be a shell of height $k$ in 
a Margulis tube or cusp in a hyperbolic 3-manifold. 
%If $T$ is a tube,
%then we assume that the distance of the shell from the core curve is at least one.
Let $f:N\to \mathbb{R}$ be a function which is smooth up to
and including the boundary. 
If $\int_{V\times [0,1]}f^2\geq 3 \int_{V\times [k-1,k]}f^2$
then ${\cal R}(f)\geq \frac{1}{3}$.
\item Let $f$ be an eigenfunction on $\hat M_{\rm thick}$ with Neumann boundary
  conditions for an eigenvalue $\lambda<1$. Let $\tau>0$ be sufficiently
  small that
the closed tubular neighborhood of radius $\tau$ of 
the boundary $\partial \hat M_{\rm thick}$ in $\hat M_{\rm thick}$ is a
shell, diffeomorphic to $\partial \hat M_{\rm thick}\times [0,\tau]$.  
Then
$\int_{\partial \hat M_{\rm thick}\times \{\tau\}}f^2\geq 
\int_{\partial \hat M_{\rm thick}\times  \{0\}}f^2$. 
%the function $t\to 
%\int_{\partial \hat M_{\rm thick}\times \{t\}}f^2$ is 
%increasing on $[0,\tau]$.  
\end{enumerate}
\end{lemma}
\begin{proof} We begin with the proof of the first part of the lemma. Define 
  \begin{align}
    s&=\min\{t\in  [0,1]\mid \int_{V\times \{t\}}f^2
\geq \frac{1}{3}\int_{V\times [0,1]}f^2\}\text{ and }\notag\\
    u&=\max\{t\in [k-1,k]\mid
    \int_{V\times \{t\}}f^2\leq \int_{V\times [k-1,k]}f^2\}.\notag
    \end{align}
Then 
\begin{equation}
\int_{V\times [0,s]}f^2\leq \frac{1}{3}\int_{V\times [0,1]}f^2\leq 
\frac{1}{3}\int_N f^2,
% \text{ and }
%\int_{V\times [u,k]}f^2\leq \int_{V\times [1,k]}f^2.
\notag\end{equation}
in particular, by the assumption on $f$ in the first part
of the lemma, we have
\begin{align}
  \int_{V\times [s,u]}f^2& \geq \int_Nf^2-
\frac{1}{3}\int_{V\times [0,1]}f^2
  -\int_{V\times [k-1,k]}f^2\notag\\
&  \geq \int_Nf^2-
  \frac{2}{3}\int_{V\times [0,1]}f^2\geq \frac{1}{3}\int_Nf^2.
\notag\end{align}

Let us compute the derivative of the function
$a(t)=\int_{V\times \{t\}}f^2$. Recall that the radial projection 
$(x,t)\in V\times \{t\}\to (x,u)\in V\times \{u\}$ 
is a homothety which
scales the Lebesgue measure of the flat torus $\rho=t$ by 
a factor $b(t,u)$, with $\frac{d}{dt}b(t,u)\vert_{u=t}=c(t)\geq 2$
(in fact, for a cusp we have $c(t)\equiv 2$, and in the
case of a Margulis tube, if the radial distance
of $V\times \{t\}$ from the core geodesic equals $R$, then
$c(t)=(\cosh^2(R)+\sinh^2(R))/\sinh(R)\cosh(R)\geq 2$, compare the proof of Lemma
\ref{partialintegration}).
We obtain
\begin{equation}\label{areaderi}
  \frac{d}{dt}\int_{V\times \{t\}}f^2=c(t)\int_{V\times \{t\}}f^2
+2\int_{V\times \{t\}}f\nu(f)\end{equation}
where $\nu$ is the outer normal field of the hypersurface $V\times \{t\}$. 

Thus as $\int_{V\times \{u\}}f^2\leq \int_{V\times [k-1,k]}f^2
\leq \int_{V\times \{s\}}f^2$, we conclude that 
\begin{equation}\label{integral4}
  \int_{V\times \{u\}}f^2-\int_{V\times \{s\}}f^2=
\int_{V\times [s,u]} c(t)f^2+2f\nu(f)\leq 0\end{equation}
and hence 
\[\int_{V\times [s,u] }c(t)f^2\leq 2\vert \int_{V\times [s,u]} f\nu(f)\vert.\]
Using $c(t)\geq 2$ for all $t$ and $\nu(f)^2\leq \Vert \nabla f\Vert^2$ and
applying the Schwarz inequality as before,
we deduce
\begin{equation}\label{integral3} 
\int_{V\times [s,u]} f^2\leq \int_{V\times [s,u]}\Vert \nabla f\Vert^2\leq \int_N\Vert \nabla f\Vert^2.
\end{equation}
As
$\int_{V\times [s,u]}f^2\geq \frac{1}{3}\int_Nf^2$,
this provides the first statement in the lemma. 

To show the second statement in the lemma,
let $f$ be an eigenfunction on $\hat M_{\rm thick}$
with Neumann boundary conditions and eigenvalue $\lambda <1$.
Assume to the contrary that there exists a number 
$\tau>0$ such that
the tubular neighborhood of radius $\tau$ about the boundary
$\partial \hat M_{\rm thick}$ of $\hat M_{\rm thick}$ is a shell
diffeomorphic to $\partial \hat M_{\rm thick}\times [0,\tau]$
and that 
\begin{equation}\label{delta1}
\delta=\int_{\partial \hat M_{\rm thick}\times \{0\}}f^2-
\int_{\partial \hat M_{\rm thick}\times \{\tau\}}f^2>0.\end{equation}
Let $t_0\in (0,\tau]$ be the smallest number for which 
the equation (\ref{delta1}) holds true, with this number $\delta$. We then have
\[\frac{d}{dt}\int_{\partial \hat M_{\rm thick}\times \{t\}}f^2\vert_{t=t_0}\leq 0.\]
By formula (\ref{areaderi}), this
implies that $\int_{\partial \hat M_{\rm thick}\times \{t_0\}}
f\nu(f)<0.$

On the other hand, as $f$ is an eigenfunction for the
eigenvalue $\lambda$, with Neumann boundary conditions,
Green's formula yields that 
\begin{align}
\int_{\partial \hat M_{\rm thick}\times [0,t_0]}f\Delta f=
-\lambda \int_{\partial \hat M_{\rm thick}\times [0,t_0]}f^2\notag \\=
\int_{\partial \hat M_{\rm thick}\times \{t_0\}}f\nu(f)-
\int_{\partial \hat M_{\rm thick}\times [0,t_0]}\Vert \nabla f\Vert^2.\notag
\end{align}
Thus $\int_{\partial \hat M_{\rm thick}\times [0,t_0]}\Vert \nabla f\Vert^2\leq 
\lambda \int_{\partial \hat M_{\rm thick}\times [0,t_0]}f^2$.
As $\lambda <1$ by assumption, this violates the 
first inequality in the formula (\ref{integral3}). Note that this
inequality applies since by our setup, the estimate in
inequality (\ref{integral4}) holds true for $u=t_0$ and $s=0$. 
\end{proof}

\begin{corollary}\label{volume2}
Let $M$ be a hyperbolic 3-manifold and let $\hat M_{\rm thick}$ be the
tubular neighborhood of radius $96$ about $M_{\rm thick}$. 
Let $N$ be the tubular neighborhood of radius one about $\partial \hat M_{\rm thick}$
in $\hat M_{\rm thick}$. 
Let $f:\hat M_{\rm thick}\to \mathbb{R}$ be a smooth function with 
${\cal R}(f)\leq 1/96$; then
%be an eigenfunction for $\Delta$ on $\hat M_{\rm thick}$ with respect to the
%eigenvalue $\lambda<1/96$, with Neumann boundary conditions. Let $N$
%be the tubular neighborhood of radius one about $\partial \hat M_{\rm thick}$
%in $\hat M_{\rm thick}$; then 
\[\int_Nf^2\leq \frac{1}{32}\int_{\hat M_{\rm thick}}f^2.\]
\end{corollary}
\begin{proof}
Parameterize $\hat M_{\rm thick}-M_{\rm thick}=W$ as $W=V\times [0,96]$
where $\partial \hat M_{\rm thick}=V\times \{0\}$. With this notation, we have
$N=V\times [0,1]$. 

We distinguish two cases. In the first case, there exists some $k\in [0,95]$
so that $\int_{V\times [0,1]}f^2\geq 3\int_{V\times [k,k+1]}f^2$. By the first
part of Lemma \ref{volumecontrol}, we have
${\cal R}(f\vert_{V\times [0,k+1]})\geq 1/3$. This implies that
\[\int_{\hat M_{\rm thick}}\Vert \nabla f\Vert^2
\geq \int_{V\times [0,k+1]} \Vert \nabla f\Vert^2\geq 
\frac{1}{3}
\int_{V\times [0,k+1]}f^2.\]

But ${\cal R}(f)\leq 1/96$, that is,
\[\int_{\hat M_{\rm thick}}\Vert \nabla f\Vert^2\leq \frac{1}{96}\int_{\hat M_{\rm thick}}f^2.\]
These two estimates together yield
\[\int_{V\times [0,1]}f^2\leq \int_{V\times [0,k+1]}f^2\leq \frac{1}{32}\int_{\hat M_{\rm thick}}f^2\]
as claimed.

In the second case, we have $\int_{V\times [0,1]}f^2\leq 3\int_{V\times [k,k+1]}f^2$ for all 
$k\leq 95$. But this implies as before that 
\[\int_{V\times [0,1]}f^2\leq \frac{1}{32}
\int_{V\times [0,96]}f^2\leq \frac{1}{32}\int_{\hat M_{\rm thick}}f^2\]
which is what we wanted to show.
\end{proof}



As a consequence, we obtain Theorem \ref{thm2} from the introduction.

\begin{proposition}\label{eigenvalue}
There is a constant $q_4>0$ with the following property.
Let $M$ be a finite volume oriented hyperbolic $3$-manifold; then for a suitable
choice of a Margulis constant, we have 
\[\lambda_k(M)\leq 
(3+q_4(\log {\rm vol}(M_{\rm thin}))\lambda_k(M_{\rm thick})\] 
for all $k$ so that $\lambda_k(M_{\rm thick})<1/96$. 
\end{proposition}
\begin{proof} 
Let $M$ be a finite volume oriented hyperbolic $3$-manifold. 
Denote by
$\hat M_{\rm thick}$ the tubular neighborhood of radius $96$ 
of the thick part
$M_{\rm thick}$ of $M$ for some choice of Margulis constant. 

Let $k\geq 1$ be such that $\lambda_k(\hat M_{\rm thick})=\lambda_k<1/96$
and let 
$f:\hat M_{\rm thick}\to \mathbb{R}$ be an eigenfunction
for the eigenvalue $\lambda_k$, with 
Neumann boundary conditions. 
Then $f$ is a smooth function on $\hat M_{\rm thick}$ 
which solves the Laplace equation 
\[\Delta(f)+
\lambda_k f=0.\] 


Our goal is to extend the function $f$ on $\hat M_{\rm thick}$ 
to a function $F$ on $M$ which is contained
in the Sobolev space ${\cal H}(M)$ of square integrable
functions on $M$, with square integrable weak derivative, in 
such a way that the Rayleigh quotient of $F$ is controlled
by the Rayleigh quotient of $f$ and hence by $\lambda_k$. 

%If $\int_{\partial \hat M_{\rm thick}}f^2=0$, then the function 
%$f$ satisfies Dirichlet boundary conditions, and let $F$ be the
%extension of $f$ which vanishes identically on 
%$M-\hat M_{\rm thick}$.  

%Assume from now on that 
%$\int_{\partial \hat M_{\rm thick}}f^2\not=0$. 

%For $\delta\in (0,3)$ let $N_\delta$ be the 
The tubular neighborhood $N$ of 
$\partial \hat M_{\rm thick}$ in $\hat M_{\rm thick}$ of radius $1$
is diffeomorphic 
%Then $N_\delta$ is 
%diffeomorphic 
to $\partial \hat M_{\rm thick}\times [0,1]$,
where the real parameter is the distance $\rho$
from the boundary 
$\partial \hat M_{\rm thick}$ of $\hat M_{\rm thick}$.  
The metric on $N$ is a warped product metric.
By Corollary \ref{volume2}, we have
\begin{equation}\label{corollary36}
\int_Nf^2\leq \frac{1}{32}\int_{\hat M_{\rm thick}}f^2.
\end{equation}

% and for every 
%$\delta <3$, the subset 
%$N_\delta= \partial \hat M_{\rm thick}\times [0,\delta]$ is the tubular
%neighborhood of $\partial \hat M_{\rm thick}$ of radius $\delta$.




Let $m>0$ be such that 
$\int_N\Vert \nabla f\Vert^2-
m\lambda_k \int_Nf^2=0$. Then we can find a number
$\delta\in [0,1]$ so that
\begin{equation}\label{delta}
\int_{\partial \hat M_{\rm thick}\times \{\delta\}}
\Vert \nabla f\Vert^2
\leq m\lambda_k\int_{\partial \hat M_{\rm thick}\times \{\delta\}}f^2.
\end{equation}
%Furthermore, using once more the
%second part of Lemma \ref{volumecontrol}, we have 
%Assume for the moment that 
%\begin{equation}\label{assumption}
%  \int_{\partial \hat M_{\rm thick}\times \{\delta\}}f^2\leq 
% \frac{1}{32} \int_{\hat M_{\rm thick}}f^2
 % =\frac{1}{m\lambda_k}\int_N\Vert \nabla f\Vert^2.
%\end{equation}

Let $V=\hat M_{\rm thick}-\partial \hat M_{\rm thick}\times [0,\delta]$.
Then $V$ is a smooth submanifold of $\hat M_{\rm thick}$ with 
boundary $\partial V=
\partial \hat M_{\rm thick}\times \{\delta\}$.
Moreover, $M-V$ is a union of Margulis tubes and cusps as before.
Extend the restriction of $f$ to $\partial V$ as in 
Proposition \ref{tube}. This yields a function 
$F:M\to (0,\infty)$ which is continuous, smooth away from 
$\partial V$, with smooth restriction to $\partial V$.
Furthermore, $F$ is square integrable, and its 
derivative (which exists in the strong sense away from the
compact hypersurface $\partial V$) is pointwise bounded. 
Thus $F$ is contained in the Sobolev space ${\cal H}(M)$ of 
square integrable functions with square integrable 
weak derivative. 


The Rayleigh quotient of $F$ can be estimated as follows. 
%Recall the definition of the number $m>0$. 
Let \[{\cal T}=(M-\hat M_{\rm thick})\cup N\supset M-V.\] 
%By the first part of Proposition
%\ref{tube} and the assumption (\ref{assumption}),  
%we have 
%\begin{equation}\int_{\cal T}F^2=\int_{{\cal T}\cap V}F^2+\int_{{\cal T}- V}F^2\leq
%  \int_Nf^2+ \frac{1}{2}\int_{\partial \hat M_{\rm thick}\times \{\delta\}}
% f^2\leq
%  \frac{3}{2}\int_Nf^2.\notag\end{equation}
%and also $\int_{\cal T}F^2\geq \frac{1}{4}
%\int_{\partial \hat M_{\rm thick}\times \{\delta\}} f^2$.
Using the second part of Proposition \ref{tube}, the estimate
(\ref{delta}) and decomposing
integrals, we have 
\begin{align}\label{formula1}
\int_{\cal T}\Vert \nabla F\Vert^2&\leq \int_{{\cal T}\cap V}\Vert \nabla f\Vert^2 +
 m\lambda_k q_2(\log {\rm vol}({\cal T}))\int_{{\cal T}-V}F^2\notag \\ 
   &\leq  m\lambda_k (1+\frac{1}{2}q_2
       \log {\rm vol}({\cal T}))
             \int_N f^2\end{align}
where the last inequality uses the definition of the constant $m>0$. 

Recall that the Rayleigh quotient of the function $f$ on 
$\hat M_{\rm thick}$ equals $\lambda_k$.
We compute
\begin{align}\label{step5}
  \int_M\Vert \nabla F\Vert^2=&
        \int_{M-{\cal T}}\Vert \nabla F\Vert^2+\int_{{\cal T}}
\Vert \nabla F\Vert^2
 \leq \int_{\hat M_{\rm thick}}\Vert \nabla f\Vert^2+
\int_{\cal T}\Vert \nabla F\Vert^2\\
&\leq 
\lambda_k\bigl(\int_{\hat M_{\rm thick}}f^2 +
m(1+\frac{1}{2}q_2\log  {\rm vol}({\cal T}))
                                    \int_Nf^2\bigr)
                                    \notag
\end{align}
where the last inequality follows from (\ref{formula1}). 
%This shows that the Rayleigh quotient of 
%$F$ is bounded from above by 
%$\lambda_k(1+m(q_2\log {\rm vol}({\cal T})+1))$. 
%If $m\leq 1$ then this shows the estimate we are aiming at.

%Assume now that $m>1$. 
As
$\int_N\Vert \nabla f\Vert^2=m\lambda_k \int_N f^2$, we have
$\int_Nf^2\leq \frac{1}{m}\int_{\hat M_{\rm thick}} f^2$. 
Inserting into the estimate (\ref{step5}) 
implies that 
\[\int_{M}\Vert \nabla F\Vert^2\leq 
\lambda_k(2+\frac{1}{2}q_2\log {\rm vol}({\cal T}))\int_{\hat M_{\rm thick}}f^2.\]

To complete the control on the Rayleigh quotient of $F$, we compare
the square norm of $F$ to the square norm of $f$. For convenience, 
extend the eigenfunction $f$ on $\hat M_{\rm thick}$ to all of $M$ by $0$.
As $F=f$ on $V=\hat M_{\rm thick}-\partial \hat M_{\rm thick}\times [0,\delta]$,
using inequality (\ref{corollary36}) we estimate
\begin{equation}\label{normcompare}
  \frac{31}{32}\int_{\hat M_{\rm thick}}f^2 \leq
  \int_{V}f^2
    \leq \int_Vf^2+\int_{M-V}F^2=\int_MF^2.
% \notag\\
% &\leq \int_Vf^2+\frac{1}{2}
%  \int_{\partial \hat M_{\rm thick}\times \{\delta\}}f^2  \leq
% \frac{65}{64}\int_Mf^2.
  \end{equation}
We deduce that 
the Rayleigh 
quotient of $F$ is bounded from above by  
\[(3+q_2(\log {\rm vol}(M_{\rm thin}))\lambda_k(\hat M_{\rm thick}).\]
This is the property of the extension function $F$ we were aiming at.
Note also for later reference that as another consequence of 
Lemma \ref{volumecontrol} and the fact that $F=f$ on $V$, we obtain
\begin{equation}\label{squarenorm2}
\int_M(f-F)^2=\int_{M-V}F^2\leq \frac{1}{2}\int_{\partial 
\hat M_{{\rm thick}\times \{\delta\}}}f^2.
%\leq \frac{1}{32}\int_Vf^2\leq 
%\frac{1}{32}\int F^2.
\end{equation}


To summarize, for an eigenfunction $f_k$ on $\hat M_{\rm thick}$
with Neumann boundary conditions and eigenvalue
$\lambda_k<1/96$, we constructed an extension $F_k:M\to \mathbb{R}$
with controlled Rayleigh quotient as predicted in the proposition.


Now let us assume that for each $i\leq k$ we constructed from an 
eigenfunction $f_i$ on $\hat M_{\rm thick}$ with eigenvalue $\lambda_i(\hat M_{\rm thick})$
with the above procedure the function $F_i$. 
Let $E_{k-1}$ be the linear span of the functions
$F_i$ for $i<k$. Assume that the function $F_k$ is normalized 
so that $\int F_k^2=1$. Using the above notations with $\delta_i\in [0,1]$
the number which enters the construction of $F_i$, 
assume furthermore for the moment
that 
\begin{equation}\label{integralestimate}
\int_{\partial \hat M_{\rm thick}\times \{\delta_k\}}f_k^2\leq 
\int_Nf_k^2\leq \frac{1}{32}\int_{\hat M_{\rm thick}}f_k^2.
\end{equation}
Using inequalities (\ref{normcompare}) and (\ref{squarenorm2}), we then
have 
\[\int_{\cal T}F_k^2\leq \frac{3}{32}\int_M F_k^2.\]
%for all $i\leq k$. 
%Using the inequality $(a-b)^2\leq 2a^2+2b^2$ we then have 
%\begin{equation}\label{squarenorm10}
%\int_M(F_i-f_i)^2=\int_{M-V}(F_i-f_i)^2\leq
%2\int_Nf_i^2+2\int_{\cal T}F_i^2\leq \frac{4}{32}.
%\end{equation}

The function $F_k$ may not be orthogonal to $E_{k-1}$ for the 
$L^2$-inner product. Denote by $H$ the $L^2$-orthogonal 
projection of $F_k$ to $E_{k-1}$.
Note that $\int_M H^2<1$. 
%To this end 
%observe first that $\int_Mf_k^2\leq \frac{32}{31}$. 
%Namely, in the case 
%that the estimate 
%(\ref{assumption}) is valid, we have
%\[\int f^2=\int_{\hat M_{\rm thick}-N}f^2+\int_Nf^2\leq F^2+\int_Nf^2\]
%and hence $\frac{31}{32}\int f^2\leq \int F^2$ by the estimate on the integral
%$\int_Nf^2$ 
%In second variation of our construction, we have
%$\int_f^2\leq \int F^2$ by immediate construction.
%
Since $H\in E_{k-1}$ there exists a finite linear combination 
$h$
of eigenfunctions on $\hat M_{\rm thick}$ with
Neumann boundary conditions and eigenvalues $\lambda_i(\hat M_{\rm thick})$ 
for $i<k$ 
which gives rise to $H$
with the above construction. As $\int_M H^2<1$, by construction 
 $\int_{\hat M_{\rm thick}} h^2\leq \frac{32}{31}$, and 
Corollary \ref{volume2} shows that
\[\int_Nh^2\leq \frac{1}{32}\int_{\hat M_{\rm thick}}h^2
\leq \frac{1}{31}.\]

Now $F_k=f_k$ and $H=h$ on $\hat M_{\rm thick}-N$ and furthermore
$\int_{\hat M_{\rm thick}}f_kh=0$ and hence
\begin{align}\label{computation10}
\int_M (F_k-H)^2 &\geq \int_{\hat M_{\rm thick}-N}(F_k-H)^2
=\int_{\hat M_{\rm thick}-N}f_k^2-2f_kh+h^2\\
& \geq \frac{29}{32} -2 \vert \int_N f_kh \vert \geq 
\frac{29}{32}-\frac{2}{31} \geq \frac{26}{32}\notag
\end{align}
by the Cauchy Schwarz inequality. 

By the above construction, there exists a 
universal constant $c>0$ so that 
\[\int_M\Vert \nabla F_k\Vert^2\leq \lambda_k(3+
c\log {\rm vol}(M_{\rm thin}))\int_{\hat M_{\rm thick}} F_k^2\] 
and the same holds true for $H$, with $\lambda_k$ replaced by
$\lambda_{k-1}$. But this implies that 
the Rayleigh quotient of $F_k-H$ is bounded
from above by a fixed multiple of $\lambda_k(M_{\rm thick})(1+\log{\rm vol}(M_{\rm thin}))$.
As $k$ with $\lambda_k<1/96$ was arbitrary, 
the proposition now follows from Rayleigh`s principle
(see \cite{C84} for details) provided that we can assure that
inequality (\ref{integralestimate}) holds true for all $k$.

%
%
%The function $F$ may not be of zero mean.
%But
%we have $\int_{M-M_{\rm thick}}F^2\leq a\int_{M_{\rm thick}}f^2$
%for a universal constant $a>0$. Moreover, 
%As $f_k$ is an eigenfunction for the eigenvalue $\lambda_k(\hat M_{\rm thick})$, 
%we have 
% and therefore
%\[\vert \int F\vert =\vert \int_{M-M_{\rm thick}}F\vert 
%\leq a\int f^2.\]
%$\int f_kh=0$.
%Since $\int F_k^2=1$, inequality (\ref{normcompare}) shows that  
%  $\int_Mf_k^2\leq
 %  \frac{32}{31}$.
%  Furthermore, $\int_M(F_k-f_k)^2\leq \frac{1}{32}$  
%  by 
%the estimate (\ref{squarenorm2}) and the assumption (\ref{integralestimate}),  
%and consequently 
%\begin{align}\label{innerproduct}
%\int_M F_kH&=\int_M (F_k-f_k+f_k)H\\
%&=\int_M (F_k-f_k)H +\int_M f_k(H-h)+f_kh \notag\\
%         &\leq (\int_M (F_k-f_k)^2)^{1/2} +\sqrt{\frac{32}{31}}
 %          (\int_M (H-h)^2)^{1/2}\leq 
%2/5\notag
%\end{align} 
%(where the last inequality follows from the estimate
%$\sqrt{\frac{1}{32}}+\sqrt{\frac{32}{31}}\sqrt{\frac{1}{32}}\leq
%\leq 2/5$).



%Now $H$ is the orthogonal projection of $F_k$ to $E_{k-1}$ 
%and therefore we have
%$\int_M H(F_k-H)=0$. Together with the estimate
%(\ref{innerproduct}), we find
%\[\int_M(F_k-H)^2=\int_MF_k(F_k-H)= 1-\int_M F_kH\geq \frac{3}{5}.\]

The final step of this proof consists in modifying the construction of the function $F$
from an eigenfunction $f$ on $\hat M_{\rm thick}$ with Neumann boundary conditions
so that $F$ fulfills inequality (\ref{integralestimate}). 


To this end recall from the second part of Lemma \ref{volumecontrol} 
that 
\[\int _{\partial \hat M_{\rm thick}\times \{0\}}f^2\leq \int_Nf^2.\]


Let $\beta:[0,\delta]\to 
\mathbb{R}$ be such that the function $u:\partial \hat M_{\rm thick}\times [0,\delta]=W
\to (0,\infty)$ defined by $u(x,s)=e^{\beta(s)}f(x,\delta)$ satisfies 
\[\int_{\partial \hat M_{\rm thick}\times \{s\}}u^2=
\int_{\partial \hat M_{\rm thick}\times \{s\}}f^2 \text{ for all }s\in [0,\delta].\]
Clearly we have $\int_Wu^2=\int_Wf^2$. 

Let $\nu$ be the vector field on $W$ which equals the 
exterior normal field of the hypersurfaces $\hat M_{\rm thick}\times \{s\}$, defining the 
orientation of the interval. We claim that 
\begin{equation}\label{w}
\int_W \nu(u)^2\leq \int_W\nu(f)^2.\end{equation}

To show the claim we evoke formula (\ref{areaderi}) from the proof of
Lemma \ref{volumecontrol} which gives

\[\frac{d}{dt}\int_{\partial \hat M_{\rm thick}\times \{t\}}f^2=
c(t)\int_{\partial \hat M_{\rm thick}\times \{t\}}f^2+2\int_{\partial \hat M_{\rm thick}\times \{t\}}f\nu(f).\]
By the definition of the function $u$, for all $t\in [0,\delta]$ we therefore have
\[\int_{\partial \hat M_{\rm thick}\times \{t\}}f\nu(f)=\int_{\partial \hat M_{\rm thick}\times \{t\}}u\nu(u).\]

But $\nu(u)(x,t)=\beta^\prime(t)u(x,t)$ and hence
\[\int_{\partial \hat M_{\rm thick}\times \{t\}}u\nu(u)=\beta^\prime(t) 
\int_{\partial \hat M_{\rm thick}\times \{t\}}u^2,\]
furthermore 
\begin{align}\label{estimate10}
\int_{\partial \hat M_{\rm thick}\times \{t\}}\nu(u)^2&=
\vert \beta^\prime(t)\vert^2\int_{\partial \hat M_{\rm thick}\times \{t\}}u^2\notag\\
&=(\int_{\partial \hat M_{\rm thick}\times \{t\}}u\nu(u))^2/
\int_{\partial \hat M_{\rm thick}\times \{t\}}u^2.\end{align}

Observe that the last line in equation (\ref{estimate10})
coincides with the corresponding 
expression for the function $f$. By the Schwarz' inequality, 
\[\vert \int_{\partial \hat M_{\rm thick}\times \{t\}}
  f\nu(f)\vert \leq
  (\int_{\partial \hat M_{\rm thick}\times \{t\}} f^2)^{1/2}
  (\int_{\partial \hat M_{\rm thick}\times \{t\}} \nu(f)^2)^{1/2}.\]
On the other hand, by (\ref{estimate10}), for the
function $u$ equality holds in this inequality. This yields 
$\int_W \nu(u)^2\leq \int_W \nu(f)^2$ as claimed. 

Following the discussion in the proof of Proposition \ref{tube}, for 
$s,t\in [0,\delta]$ the radial 
projection of the hypersurface
$\partial \hat M_{\rm thick}\times \{s\}$ onto the 
hypersurface 
$\partial \hat M_{\rm thick}\times \{t\}$ is bilipschitz, with uniformly bounded 
bilipschitz constant not depending on $M$. 
This implies that there is a universal constant
$c>0$ such that for each $s\in [0,\delta]$, 
the Rayleigh quotient
of the restriction of $u$ to
$\partial \hat M_{\rm thick}\times \{s\}$ is not bigger than
$c$ times the Rayleigh quotient of the restriction of $u$ to 
$\partial \hat M_{\rm thick}\times \{\delta\}$.
Together with the estimate (\ref{w})
on normal derivatives and the choice of
$\delta$, we obtain 
\[\int_{\partial \hat M_{\rm thick}\times [0,\delta]}
  \Vert \nabla u\Vert^2
  \leq cm\lambda_k\int_{\partial \hat M_{\rm thick}\times
    [0,\delta]}f^2
  +\int_{\partial \hat M_{\rm thick}\times [0,\delta]}
  \Vert \nabla f\Vert^2.\]


Use Proposition \ref{tube} to 
extend the restriction of the
function $u$ to $\hat M_{\rm thick}\times \{0\}$ to 
the complement of 
$\hat M_{\rm thick}$ in $M$
(which consists of a collection of Margulis tubes and cusps).
It now follows from the beginning of this proof that the
resulting function $F$ has all properties
predicted in
the proposition.

To provide more details of this estimate, let 
again ${\cal T}=M-\hat M_{\rm thick}\cup N$. 
Then $f=F$ on $M-{\cal T}$, and by Proposition \ref{tube},
Lemma \ref{volumecontrol}
and inequality (\ref{corollary36}), 
\begin{equation}\label{estimate12}
  \int_{\cal T}F^2\leq \frac{3}{2}\int_Nf^2\leq
\frac{3}{64}\int_{\hat M_{\rm thick}}f^2. \end{equation}
This implies as in the estimate (\ref{normcompare}) that 
\begin{equation}\label{comparisoncontrol2}
\int_Mf^2\leq \int_MF^2\leq \frac{67}{64}\int_Mf^2.
%  \int_M(F-f)^2\leq 2\int_{\cal T}F^2+2\int_Nf^2
%  \frac{5}{32}\int_M f^2.
\end{equation}
On the other hand, using the estimate (\ref{delta})
and the construction of $F$, 
\[\int_{{\cal T}}\Vert \nabla F\Vert^2\leq
  \lambda_kq_2m(\log{\rm vol}({\cal T})+c)\int_{N}f^2\]
which yields the required estimate on Rayleigh quotients as above.
\end{proof}



\begin{remark} For a compact 
hyperbolic manifold $M$ of dimension $n\geq 4$, 
Burger and Schroeder 
\cite{BS87} proved the upper bound
\[\lambda_1(M)\leq \frac{\alpha_n+\beta_n\log{\rm vol}(M)}{{\rm diam}(M)}\]
with constants $\alpha_n,\beta_n$ depending only on $n$.
As hyperbolic $3$-manifolds with cusps can be Dehn-filled, with 
fixed volume bound,  this result is in general false for 
hyperbolic $3$-manifolds.
\end{remark}


%%%%%%%%%%%%%%%%%%%%%%%%%%%%%%%%%%%%%%%%%%%%%%%%%%%%%%%%%%%%%%%%%%%%%%%%%%%
%%%%%%%%%%%%%%%%%%%%%%%%%%%%%%%%%%%%%%%%%%%%%%%%%%%%%%%%%%%%%%%%%%%%%%%%%%%

 
 
 
 

\begin{thebibliography}{28}


\bibitem{BGH} H.~Baik, I.~Gekhtman, U.~Hamenst\"adt,
{\em The smallest positive eigenvalue of fibered
hyperbolic 3-manifolds}, Proc. Lond. Math. Soc. 120 (2020), 704--741.

%\bibitem{BP} R.~Benedetti, C.~Petronio, {\em Lectures on hyperbolic
%geometry}, Springer Universitext, Springer Berlin Heidelberg 1992.


%\bibitem{B03} J. Brock {\em Weil-Petersson translation distance and volumes of mapping tori}, Comm. Anal. Geom. 11 (2003), pp. 987-999. 
%\bibitem{BB04} J. Brock, K. Bromberg (2004). \emph{On the density of geometrica%lly finite Kleinian groups.} Acta Math. 192, pp. 33?93.

%\bibitem{BB14} J.~Brock, K.~Bromberg, {\em Inflexibility, Weil-Peterson distance, 
%and volumes of fibered 3-manifolds}, arXiv:1412.0733, to appear i Math. Res. Letters


%\bibitem{BBCL14} J. Brock, K. Bromberg, R. Canary, C. Lecuire, {\em Convergence and divergence of Kleinian surface groups}, 
%J. Topol. 8 (2015), 811--841.


%\bibitem{BBCM11} J. Brock, K. Bromberg, R. Canary, Y. Minsky, {\em Local topology in deformation spaces of hyperbolic 3-manifolds}, 
%Geom. Topol., 15(2) (2011). pp. 1169--1224. 


%\bibitem{BBCM13} J. Brock, K. Bromberg, R. Canary, Y. Minsky, {\em Convergence properties of end invariants.} Geom. Topol. 17
%(2013), pp. 2877-2922. 




%\bibitem{EL2} J. Brock, R. Canary,  Y. Minsky,  {\em The classification of Kleinian surface groups, II: The Ending Lamination Conjecture},  
%Ann. Math. 176 (1) (2012), pp. 1-149. 


\bibitem{BGLM02} M.~Burger, T.~Gelander, A.~Lubotzky, S.~Mozes,
{\em Counting hyperbolic manifolds}, Geom. Funct. Anal. 12 (2002), 1161--1173.

\bibitem{BS87} M.~Burger, V.~Schroeder, {\em Volume, diameter
and the first eigenvalue of locally symmetric
spaces of rank one}, J. Diff. Geom. 26 (1987), 273--284.


%\bibitem{Bu82} P. Buser, {\em A note on the isoperimetric constant},
%Ann. Scient. E. Norm. Sup. 15 (1982), 213--230.


\bibitem{Bu} P.~Buser, {\em Geometry and spectra
of compact Riemann surfaces}, Birkh\"auser, Boston 1992.



\bibitem{BCD93} P.~Buser, B.~Colbois, J.~Dodziuk,
{\em Tubes and eigenvalues for negatively curved manifolds},
J. Geom. Anal. 3 (1993), 1--26.  


\bibitem{C84} I.~Chavel, {\em Eigenvalues in Riemannian geometry}, 
Academic Press, Orlando 1984.  

\bibitem{CD94} I. Chavel, J. Dodziuk, 
{\em The spectrum of degenerating hyperbolic 3-manifolds}, 
J. Diff. Geom. 39 (1994), 123--137. 


\bibitem{Chu96} F.R.K. Chung, {\em Laplacians on graphs and Cheeger inequality}, 
in Combinatorics, Paul Erdos is eighty, Vol 2, Bolyai Math. Soc., Budapest 1996, 157--172.

%\bibitem{Horbez-Dahmani} F.~Dahmani, C.~Horbez,  
%{\em Spectral theorems for random walks on mapping class groups and $Out(F_N)$}. 
%arxiv:1506.06790


\bibitem{D87} J. Dodziuk, {\em A lower bound for the first eigenvalue of a finite-volume
negatively curved manifold}, Bol. Soc. Bras. Mat. 18 (1987), 23--34.

\bibitem{DR86} J. Dodziuk, B. Randol, {\em Lower bounds for $\lambda_1$ on a finite-volume hyperbolic 
manifold}, J. Diff. Geom. 24 (1986), 133--139.


\bibitem{E87} J.~Eschenburg, {\em Comparison theorems and
applications}, Manuscripta Math. 59 (1987), 295--323.

\bibitem{E94} J. Eschenburg, {\em Comparison theorems in Riemannian geometry}, 
Lecture Notes Series no. 3, Universita degli Studi di Trento 1994. 


%\bibitem{FS13} D. Futer, S. Schleimer, \emph{Cusp geometry of fibered 3-manifolds}. Amer. J. Math. 136 (2014), Issue 2, pp. 309-356. 

\bibitem{G04} T.~Gelander, {\em Homotopy
type and volume of locally symmetric manifolds},
Duke Math. J. 124 (2004), 459--515.


\bibitem{H04} U.~Hamenst\"adt, {\em Small eigenvalues of 
geometrically finite manifolds}, J. Geom. Anal. 14 (2004), 281--290. 


%\bibitem{Hamenstaedt10} U.~Hamenst\"adt,
%{\em Stability of quasi-geodesics in Teichm\"uller space},
%Geom. Dedicata 146 (2010), 101--116.

%\bibitem{Hamenstaedt13} U. Hamenst\"adt (2013). \emph{Bowen's construction for the Teichm\"uller flow.} J. Mod. Dynamics 7, pp. 489-526. 

%\bibitem{Hamenstaedt16} U.~Hamenst\"adt, 
%{\em Counting periodic orbits in the thin part of strata},
%preprint 2016.




%\bibitem{Hamenstaedt05} U. Hamenst\"adt. \emph{Word hyperbolic extensions of su%rface groups.} arXiv:math/0505244v2. 

\bibitem{HV19} U. Hamenst\"adt, G. Viaggi, {\em Small eigenvalues of random 3-manifolds},
arXiv:1903.08031.


%\bibitem{Kaimanovich} V. A. Kaimanovich, {\em The Poisson formula for groups with hyperbolic
%properties}, Ann. of Math. (2) 152 (2000), 659--692.

%\bibitem{Kaimanovich-Masur} V.A. Kaimanovich  and H. Masur, {\em 
%The Poisson boundary of the mapping class group}, Invent. Math. 125 (1996), 221--264.


%\bibitem{KMcS14} Kojima, G.~McShane, {\em Normalized entropy versus volume for pseudo-Anosovs},
%arXiv:1411.6350.
 



%\bibitem{L06} M.~Lackenby, {\em Heegaard splittings, the virtually Haken
%conjecture and Property $(\tau)$}, 
%Invent. Math. 164 (2006), 317--369.


\bibitem{LS16} A.~Lenzhen, J.~Souto, to appear.

\bibitem{LS12} L.~Louder, J.~Souto, {\em Diameter and spectral gap for planar graphs}, arXiv:1204.4435v2 [math.GT].


%\bibitem{Ma11} J.~Maher, {\em Random walks on the mapping class
%group}, Duke Math. J. 156 (2011), 429--486.

%\bibitem{Maher-Tiozzo}
%J. Maher and G. Tiozzo, {\em Random walks on weakly hyperbolic groups},
%arXiv:1410.4173v2.



\bibitem{Man05} T.~ Mantuano, {\em Discretization of Compact Riemannian Manifolds Applied to the Spectrum of Laplacian}, Ann.  Global 
Anal. \& Geom. 27 (2005), 33--46.



%\bibitem{MM99} H. Masur, Y. Minsky, \emph{Geometry of the complex of curves I: Hyperbolicity.} Invent. Math. 138 (1) (1999), 103--149. 
%\bibitem{MM00} H. Masur, Y. Minsky, \emph{Geometry of the complex of curves II: Hierarchical structure.} Geom. Funct. Anal. 10 (4) (2000), 902--974. 


%\bibitem{Minsky93} Y. Minsky, {\em Teichm\"uller geodesics and end invariants of hyperbolic 3-manifold.} Topology 32 (3) (1993), pp.625-647. 

%\bibitem{Minsky94} Y. Minsky, {\em On rigidity, limit sets, and end invariants of hyperbolic 3-manifolds}, J. Amer. Math. Soc. 7 (3) (1994).



%\bibitem{Minsky96} Y.~Minsky, {\em Quasi-projections in 
%Teichm\"uller space}, J. Reine Angew. Math. 473 (1996), 
%121--136. 



%\bibitem{EL1} Y. Minsky, {\em The classification of Kleinian surface groups, I: Models and bounds}, Ann. Math. 171 (1) (2010), pp. 1-107. 

\bibitem{OR09} J.~P.~Otal, E.~Rosas,
{\em Pour toute surface hyperbolique de genre $g$, 
$\lambda_{2g-2}>1/4$}, Duke Math. J. 150 (2009), 101--115. 


%\bibitem{PH92} R.~Penner, J.~Harer, {\em Combinatorics of train tracks},
%Ann. Math. Studies 1992.

%\bibitem{Rafi14} K. Rafi (2014). \emph{Hyperbolicity in Teichm\"uller space.} G%eom. Topol. 18, pp. 3025-3053.

%\bibitem{Rivin} I. Rivin, \emph{Statistics of Random 3-Manifolds occasionally  fibering over the circle.} arXiv:1401.5736 [math.GT]. 

\bibitem{S82} R. Schoen,  {\em A lower bound for the first eigenvalue of a negatively curved manifold}, 
J. Differential Geom. 17 (2) (1982), 233-238.


%\bibitem{Tiozzo} G. Tiozzo, 
%{\em Sublinear deviation between geodesics and sample paths},
%Duke Mathematical Journal 164 (2015), no. 3, 511-539 


\bibitem{White} N. White, {\em Spectral bounds on closed hyperbolic 3-manifolds}, J. London Math. Soc. 87 (2013), 837--853.  



%\bibitem{Rafi05}  K. Rafi (2005), \emph{A characterization of short curves of a Teichm\"uller geodesic.} Geom. Topol. 9, pp. 179-202. 





%\bibitem{Canary96} R. Canary (1996). \emph{A covering theorem for hyperbolic 3-manifolds and its applications.} Topology 35 (3), pp.751-778. 

%\bibitem{DDM} S. Dowdall, M. Duchin, H. Masur. \emph{Statistically hyperbolicity in Teichm\"uller space.} arXiv:1108.5416v5%%. 

%\bibitem{HK05} C. Hodgson, S. Kerckhoff (2005). \emph{Universal bounds for hyperbolic Dehn surgery.} Ann. Math. 162 (1), pp. 367-421.
%\bibitem{Rafi07} K. Rafi (2007), \emph{Thick-think decomposition for quadratic differentials.} Math. Res. Lett. 14, no. 2, pp. 333-341. 

\end{thebibliography}

\bigskip

\noindent
Universit\"at Bonn, Endenicher Allee 60, 53115 Bonn, Germany\\
e-mail: ursula@math.uni-bonn.de



\end{document}

%%%%%%%%%%%%%%%%%%%%%%%%%%%%%%%%%%%%%%%%%%%%%%%%%%%%%%
%%%%%%%%%%%%%%%%%%%%%%%%%%%%%%%%%%%%%%%%%%%%%%%%%%%%%%%

